\documentclass[11pt,a4paper,reqno]{amsart}

\usepackage[utf8]{inputenc}
\usepackage[T1]{fontenc}
\usepackage{amsmath,amssymb,amsfonts,amsthm,mathtools}
\usepackage{geometry}
\geometry{margin=2.5cm}
\usepackage[dvipsnames]{xcolor}
\usepackage{hyperref}
\hypersetup{
    colorlinks=true,
    linkcolor=blue!70!black,
    citecolor=green!50!black,
    urlcolor=blue!70!black
}
\usepackage{microtype}
\usepackage{booktabs}
\usepackage{array}
\usepackage{physics}

% Theorem Environments
\newtheorem{theorem}{Theorem}[section]
\newtheorem{lemma}[theorem]{Lemma}
\newtheorem{proposition}[theorem]{Proposition}
\newtheorem{corollary}[theorem]{Corollary}

\theoremstyle{definition}
\newtheorem{definition}[theorem]{Definition}
\newtheorem{example}[theorem]{Example}

\theoremstyle{remark}
\newtheorem{remark}[theorem]{Remark}

\numberwithin{equation}{section}
\raggedbottom

% Custom Operators
\newcommand{\dif}{\mathrm{d}}
\DeclareMathOperator{\Sym}{Sym}
\DeclareMathOperator{\End}{End}
\DeclareMathOperator{\supp}{supp}
\DeclareMathOperator{\im}{im}
\DeclareMathOperator{\Cov}{Cov}
\DeclareMathOperator{\Var}{Var}
\DeclareMathOperator{\Unp}{Unp}
\DeclareMathOperator{\Proj}{Proj}
\DeclareMathOperator{\reach}{reach}

\title[Beyond the Spectrum III]{Beyond the Spectrum III: Higher Invariants, Bipartite Kernels, and Non-Equilibrium Geometry of Functional Realizations}

\author{Reinaldo M. Silva-Filho}
\address{Master's student, Postgraduate Program in Statistics and Agricultural Experimentation (PPGEE/DES), Department of Statistics (DES), Federal University of Lavras (UFLA), Lavras, MG, Brazil}
\email{reinaldo.filho1@estudante.ufla.br}

\date{}

\begin{document}

\begin{abstract}
This third volume studies functional realizations of matrices through symplectic, microlocal, nonlinear spectral, stochastic, bipartite-kernel and geometric-measure invariants. The main statements, proved here or derived from cited classical results, are these. The first Ekeland--Hofer capacity of a convex sublevel set $K_t(A)\subset\mathbb{R}^{2n}$ lies between $2\pi(t-\min\Phi(A))/\lambda_{\max}$ and $2\pi(t-\min\Phi(A))/\lambda_{\min}$, where $\lambda_{\min},\lambda_{\max}$ bound the Hessian on $K_t(A)$. Rank stratifications of matrix-valued realizations transverse to the rank stratification of $\mathrm{Sym}_n(\mathbb{R})$ are Whitney regular. For a weighted $p$-Laplacian whose weight is bounded above and below, $(\lambda_1^{(p)})^{1/p}$ converges as $p\to\infty$ to the reciprocal inradius of the domain, whatever the weight; on the unit disk this limit is $1$, while the Cheeger constant is $2$. For $\kappa$-convex Langevin potentials on $\mathbb{R}^d$ the law converges in $\mathcal{W}_2$ at rate $\kappa$, and the Jarzynski identity gives $e^{-\Delta F}=\Tr(B)/\Tr(A)$ for trace-normalized realizations. For bipartite kernels, the partial trace has trace $\Tr(A)$ and rank at most the sum of the Schmidt ranks of the frame, and the normalized modular Hamiltonian is non-negative. The simplicial Mellin transform $\mathcal{Z}_A(s)$ is entire. Compact regular level sets have reach at least $\min(\epsilon_0/M, d_{\mathrm{sep}}/2)$. Counterexamples show that three stronger statements fail: a stratified Euler-characteristic formula, a thermodynamic-length bound carrying a curvature factor, and a finite-rank bound for partial traces of general frames. The holographic relation $S_R=2E_W$ is recalled as a physical proposal, not proved.
\end{abstract}

\maketitle
\tableofcontents

% =========================================================================
% SECTION 1
% =========================================================================
\section{Introduction}
\label{sec:intro}

Throughout, $\mathcal{S}_+^n$ denotes the cone of positive semi-definite Hermitian $n\times n$ matrices. The classical spectral theory of matrices and tensors characterizes an algebraic object $A \in \mathcal{S}_+^n$ through scalar invariants invariant under orthogonal or unitary transformations: eigenvalues $\lambda_j(A)$, traces $\Tr(A^k)$, and determinants $\det(A)$. In Volumes~I and~II of this series \cite{SilvaFilhoBTS}, we introduced the paradigm of \emph{functional realizations}, embedding finite- and infinite-dimensional tensors into continuous function spaces:
\begin{equation}
\Phi \colon \mathcal{T}^k(V) \longrightarrow \mathcal{F}(\Omega, \mathbb{R}), \qquad A \longmapsto \Phi(A)(x),
\end{equation}
where $\Omega \subset \mathbb{R}^d$ is a smooth compact domain or Riemannian manifold. This continuous lifting unlocked the rich geometric measures of Sobolev spaces $W^{1,p}$, bounded variation $\mathrm{BV}(\Omega)$, continuous optimal transport $\mathcal{W}_2(\mu_A, \mu_B)$, Bakry--Émery Ricci curvature lower bounds $\kappa_{\mathrm{BE}}(A)$, persistent homology diagrams $\mathrm{Dgm}_k(A)$, integrated Willmore bending energies $\mathcal{W}(\Phi(A))$, and non-commutative Dixmier traces $\Tr_\omega(\Phi(A) |\mathcal{D}|^{-d})$.

Despite the analytical breadth of these invariants, Volumes~I and~II operated under two fundamental structural constraints:
\begin{enumerate}
    \item \textbf{Mostly Scalar Codomains:} The realization $\Phi(A)$ was mostly a scalar field $\Omega \to \mathbb{R}$ (Volume~I also uses integral operators with kernel $W_A$) or a Radon probability measure $\dif\mu_A = \Phi(A) \dif x$. Consequently, topological phase-space rigidity (symplectic capacity), gauge-field parallel transport, and internal quantum entanglement remained inaccessible.
    \item \textbf{Static Equilibrium Assumption:} Invariants were computed on static tensors or quasistatic limits, blind to finite-time non-equilibrium dissipation, dynamical fluctuations, and irreversible entropy production.
\end{enumerate}

The present volume (\emph{Beyond the Spectrum III}) addresses these two constraints along three fronts:
\begin{itemize}
    \item \textbf{Symplectic invariants:} In Section~\ref{sec:symplectic}, we treat $\Omega = \mathbb{R}^{2n}$ as a symplectic phase space $(\Omega, \omega_0)$ and bound the first Ekeland--Hofer capacity $c_1^{\mathrm{EH}}$ of the sublevel sets $K_t(A)$ by Hessian data; we recall the Hofer--Zehnder capacity of convex bodies.
    \item \textbf{Matrix-valued realizations and bipartite kernels:} In Section~\ref{sec:microlocal}, $\Phi(A)(x)$ is a smooth matrix-valued field in $\mathrm{Sym}_n(\mathbb{R})$, giving Whitney rank stratifications and constructible sheaves with Lagrangian micro-support. In Section~\ref{sec:bipartite}, we pass to bipartite integral kernels $K_A(x, y)$, their partial traces, modular Hamiltonians and reflected entropies; the holographic relation with the entanglement wedge cross-section is recalled as a physical proposal.
    \item \textbf{Nonlinear analysis, stochastic thermodynamics, and geometric measure theory:} In Section~\ref{sec:nonlinear}, we study the weighted $p$-Laplacian family $\{\lambda_1^{(p)}(A)\}_{p \in (1,\infty)}$ and its limit as $p\to\infty$, and Gromov $\delta$-hyperbolicity. In Section~\ref{sec:thermo}, we consider Langevin diffusions under the potential $V_A = -\log \Phi(A)$, Jarzynski's identity, and the slow-driving limit of the dissipated work. Finally, Sections~\ref{sec:simplicial} and~\ref{sec:federer} treat Mellin transforms over the simplex and the reach of level sets.
\end{itemize}

Statements carry tags of the form $\mathtt{OBL\text{-}nnn}$ for cross-reference. Results labelled Theorem, Proposition or Lemma are proved in the text or are cited classical results with the source named; statements that fail in general are recorded as remarks together with a counterexample.

\subsection*{Related work}
Functional realizations are also the subject of Chapter~1 of the author's monograph~\cite{silvafilho2026book}. Most results below are classical or follow from classical results, and we use them rather than re-derive them.

Symplectic capacities originate in Gromov's non-squeezing theorem~\cite{Gromov1985}; the capacities used here are those of Ekeland and Hofer~\cite{EkelandHofer1989,EkelandHofer1990} and of Hofer and Zehnder~\cite{HoferZehnder1994}, surveyed by Cieliebak, Hofer, Latschev and Schlenk~\cite{CieliebakHoferLatschevSchlenk2005}. On convex domains with smooth boundary the first Ekeland--Hofer capacity and the Hofer--Zehnder capacity coincide and equal the minimal action of a closed characteristic on the boundary \cite[Theorem~1.3]{ArtsteinAvidanOstrover2008}. For convex bodies, Artstein-Avidan and Ostrover proved a Brunn--Minkowski inequality for this capacity~\cite{ArtsteinAvidanOstrover2008}, and Artstein-Avidan, Karasev and Ostrover related its maximization at fixed volume to Mahler's conjecture~\cite{ArtsteinAvidanKarasevOstrover2014}. Viterbo~\cite{Viterbo2000} suggested that among convex domains of given volume the ball maximizes capacity; Haim-Kislev and Ostrover~\cite{HaimKislevOstrover2026} gave a counterexample to Viterbo's volume--capacity conjecture.

Whitney's conditions~(a) and~(b) go back to Whitney~\cite{Whitney1965}; see Mather's notes~\cite{Mather2012}, the Liverpool notes~\cite{GibsonEtAl1976} and Trotman's survey~\cite{Trotman2020}. Micro-supports, constructible sheaves and characteristic cycles are developed by Kashiwara and Schapira~\cite{KashiwaraSchapira1990}; the index theorem for constructible sheaves is due to Kashiwara~\cite{Kashiwara1985}, and Schapira~\cite{Schapira2017} gives an overview. Integration of constructible functions with respect to the Euler characteristic was developed by Viro~\cite{Viro1988} and Schapira~\cite{Schapira1991}.

The limit $p\to\infty$ of the first Dirichlet eigenvalue of the $p$-Laplacian was found by Juutinen, Lindqvist and Manfredi~\cite{JuutinenLindqvistManfredi1999} and by Fukagai, Ito and Narukawa~\cite{FukagaiItoNarukawa1999}; see Lindqvist's lecture notes~\cite{Lindqvist2008}, the Finsler version of Belloni, Kawohl and Juutinen~\cite{BelloniKawohlJuutinen2006}, and the optimal-transport approach of Champion, De~Pascale and Jimenez~\cite{ChampionDePascaleJimenez2008}. The opposite limit $p\to1$ gives the Cheeger constant~\cite{Cheeger1970,KawohlFridman2003}.

On the stochastic side, the Fokker--Planck equation of an overdamped Langevin diffusion is the $\mathcal{W}_2$-gradient flow of the relative entropy~\cite{JordanKinderlehrerOtto1998}, and lower Ricci bounds are equivalent to $\mathcal{W}_2$ contraction estimates for the heat flow on Riemannian manifolds~\cite{vonRenesseSturm2005}; see Villani's monograph~\cite{Villani2009}. Stochastic thermodynamics is reviewed by Seifert~\cite{Seifert2012}. Thermodynamic length goes back to Salamon and Berry~\cite{SalamonBerry1983}; Crooks~\cite{Crooks2007} defined it for small systems described by equilibrium statistical mechanics, and Sivak and Crooks~\cite{SivakCrooks2012} derived the friction tensor that controls dissipation in linear response. Optimal finite-time protocols were computed by Schmiedl and Seifert~\cite{SchmiedlSeifert2007}, and their link with optimal transport was found by Aurell, Mej\'ia-Monasterio and Muratore-Ginanneschi~\cite{Aurell2011}. Recent work relates entropy production, speed limits and optimal transport~\cite{NakazatoIto2021,VanVuSaito2023}, and shows that for overdamped dynamics the friction-tensor geometry of the slow-driving limit coincides with $L^2$ optimal-transport geometry~\cite{ZhongDeWeese2024}; Blaber and Sivak~\cite{BlaberSivak2023} review optimal control in this setting.

The reflected entropy was introduced by Dutta and Faulkner~\cite{DuttaFaulkner2019}. The entanglement wedge cross-section was proposed as the holographic dual of the entanglement of purification by Umemoto and Takayanagi~\cite{UmemotoTakayanagi2018} and by Nguyen, Devakul, Halbasch, Zaletel and Swingle~\cite{NguyenEtAl2018}. Akers and Rath~\cite{AkersRath2020} and Hayden, Parrikar and Sorce~\cite{HaydenParrikarSorce2021} studied the gap between reflected entropy and mutual information in holographic states, and Hayden, Lemm and Sorce~\cite{HaydenLemmSorce2023} showed by counterexample that the reflected entropy is not monotonically decreasing under partial trace. Sets of positive reach were introduced by Federer~\cite{Federer1959}; see the monograph of Rataj and Z\"ahle~\cite{RatajZahle2019}. The reach controls the curvature of a submanifold~\cite{NiyogiSmaleWeinberger2008,Aamari2019} and the variation of its tangent spaces~\cite{BoissonnatLieutierWintraecken2019}. The tube formula for compact submanifolds is due to Weyl~\cite{Weyl1939}.

% =========================================================================
% SECTION 2
% =========================================================================
\section{Symplectic Invariants of Functional Sublevel Sets}
\label{sec:symplectic}

Let $(\mathbb{R}^{2n}, \omega_0)$ be the standard symplectic vector space with canonical coordinates $x = (q_1, \dots, q_n, p_1, \dots, p_n)$ and non-degenerate closed 2-form $\omega_0 = \sum_{j=1}^n \dif p_j \wedge \dif q_j$. 
Let $A \in \mathcal{S}_+^n$, and let $\Phi(A) \in C^2(\mathbb{R}^{2n}, \mathbb{R})$ be a strictly convex functional realization satisfying $\nabla^2 \Phi(A)(x) \ge \mu I > 0$ for all $x \in \mathbb{R}^{2n}$, with $\lim_{\|x\| \to \infty} \Phi(A)(x) = +\infty$.

The minimum is attained at a unique point $x_0$, and $\nabla\Phi(A)$ vanishes only there, so every $t > \min_{x} \Phi(A)(x)$ is a regular value. The sublevel set
\begin{equation}
K_t(A) \coloneqq \left\{ x \in \mathbb{R}^{2n} \mid \Phi(A)(x) \le t \right\}
\end{equation}
is a compact, strictly convex body in $\mathbb{R}^{2n}$ with $C^2$ boundary $\Sigma_t \coloneqq \partial K_t(A)$. On this compact body we define the Hessian spectral bounds:
\begin{equation}
\lambda_{\max} \coloneqq \sup_{x \in K_t(A)} \lambda_{\max}\left(\nabla^2 \Phi(A)(x)\right) < \infty, \qquad \lambda_{\min} \coloneqq \inf_{x \in K_t(A)} \lambda_{\min}\left(\nabla^2 \Phi(A)(x)\right) \ge \mu > 0.
\end{equation}

\begin{theorem}[\textbf{OBL-001: Sublevel Convexity and Ekeland--Hofer Capacity}]
\label{thm:eh_capacity_pos}
Let $A \in \mathcal{S}_+^n$ and let $\Phi(A) \in C^2(\mathbb{R}^{2n})$ be strictly convex with $\nabla^2 \Phi(A) \ge \mu I > 0$. For every $t > \min \Phi(A)$, the first Ekeland--Hofer capacity $c_1^{\mathrm{EH}}$ of \cite{EkelandHofer1989,EkelandHofer1990} satisfies the two-sided bound
\begin{equation}
\label{eq:eh_bound}
0 < \frac{2\pi (t - \min \Phi(A))}{\lambda_{\max}} \le c_1^{\mathrm{EH}}(K_t(A)) \le \frac{2\pi (t - \min \Phi(A))}{\lambda_{\min}} < \infty.
\end{equation}
\end{theorem}

\begin{proof}
Write $m = \min\Phi(A) = \Phi(A)(x_0)$, $\Phi = \Phi(A)$, and set $r_{\min} = \sqrt{2(t - m)/\lambda_{\max}}$, $R_{\max} = \sqrt{2(t - m)/\lambda_{\min}}$. For a unit vector $v$ and $s \ge 0$ let $f(s) = \Phi(x_0 + s v)$; then $f(0) = m$, $f'(0) = 0$ and $f''(s) = \langle \nabla^2\Phi(x_0+sv)v, v\rangle$.

\emph{Inner ball.} Suppose some point $x_0 + s_1 v$ with $s_1 < r_{\min}$ lies outside $K_t(A)$. Since $f(0) < t < f(s_1)$, there is a first $s^* \in (0, s_1)$ with $f(s^*) = t$, and the segment $\{x_0 + sv : 0 \le s \le s^*\}$ lies in $K_t(A)$, where $f'' \le \lambda_{\max}$. Taylor's formula gives $t = f(s^*) \le m + \lambda_{\max}(s^*)^2/2 < m + \lambda_{\max} r_{\min}^2/2 = t$, a contradiction. Hence the open ball $B(x_0, r_{\min}) \subseteq K_t(A)$.

\emph{Outer ball.} Let $x_0 + s v \in K_t(A)$. By convexity the whole segment from $x_0$ lies in $K_t(A)$, where $f'' \ge \lambda_{\min}$, so $t \ge f(s) \ge m + \lambda_{\min} s^2/2$, that is, $s \le R_{\max}$. Hence $K_t(A) \subseteq \bar{B}(x_0, R_{\max}) \subset B(x_0, R_{\max} + \varepsilon)$ for every $\varepsilon > 0$.

The capacity $c_1^{\mathrm{EH}}$ is monotone under inclusion, invariant under symplectomorphisms (in particular translations), and satisfies $c_1^{\mathrm{EH}}(B^{2n}(r)) = \pi r^2$ \cite{EkelandHofer1989}. Therefore $\pi r_{\min}^2 \le c_1^{\mathrm{EH}}(K_t(A)) \le \pi (R_{\max} + \varepsilon)^2$ for every $\varepsilon > 0$, which is \eqref{eq:eh_bound}.
\end{proof}

\begin{remark}
The inclusions in the proof are between round balls; the Hessian bounds do not give ellipsoids adapted to $\Phi(A)$. For a quadratic $\Phi(A)(x) = \tfrac12\langle Qx, x\rangle$ the bound is not sharp in general, since $c_1^{\mathrm{EH}}(K_t)$ is then $2\pi t$ divided by the largest symplectic eigenvalue of $Q$, which lies between $\lambda_{\min}(Q)$ and $\lambda_{\max}(Q)$.
\end{remark}

\begin{proposition}[\textbf{OBL-002: Hofer--Zehnder Periodic Orbit Invariance}]
\label{prop:hz_symp_inv}
Assume in addition that $\Phi(A)$ is $C^\infty$, and let $\psi \in \mathrm{Symp}(\mathbb{R}^{2n}, \omega_0)$ be any smooth symplectomorphism ($\psi^* \omega_0 = \omega_0$). Then the Hofer--Zehnder capacity $c_{\mathrm{HZ}}$ of the sublevel set satisfies
\begin{equation}
c_{\mathrm{HZ}}(\psi(K_t(A))) = c_{\mathrm{HZ}}(K_t(A)) = \inf_{\gamma \subset \partial K_t(A)} \left| \oint_\gamma p \dif q \right|,
\end{equation}
where the infimum runs over all closed characteristic curves of the Hamiltonian vector field $X_{\Phi(A)}$ on $\partial K_t(A)$.
\end{proposition}

\begin{proof}
This is a cited classical result. The first equality is the symplectic invariance axiom of the Hofer--Zehnder capacity. The second is the theorem of Hofer and Zehnder \cite{HoferZehnder1994} that, for a convex body with smooth boundary in $\mathbb{R}^{2n}$, $c_{\mathrm{HZ}}$ equals the minimal action of a closed characteristic on the boundary; closed characteristics on the regular level set $\partial K_t(A)$ are the periodic orbits of $X_{\Phi(A)}$ there.
\end{proof}

\begin{remark}[\textbf{OBL-003: spectral invariants of realization Hamiltonians}]
\label{rem:viterbo_lipschitz}
One would like a $C^0$-Lipschitz estimate $\lvert c(a, \Phi(A)) - c(a, \Phi(B))\rvert \le \lVert\Phi(A) - \Phi(B)\rVert$ for spectral invariants of the Hamiltonians $\Phi(A)$, $\Phi(B)$ with quadratic growth at infinity. We do not prove such an estimate here. For spectral invariants defined through generating functions or Floer homology see Viterbo~\cite{Viterbo1992} and Schwarz~\cite{Schwarz2000}, and for the symplectic homology of open sets in $\mathbb{C}^n$ see Floer and Hofer~\cite{FloerHofer1994}. A precise statement needs at least: (i) a class of Hamiltonians for which the spectral invariants $c(a,\cdot)$ are defined, for example asymptotically quadratic Hamiltonians with a fixed non-resonant quadratic form at infinity, since the filtered Floer homology changes when the slope at infinity crosses a resonance; (ii) a homology class $a$ of the relevant filtered Floer complex (the singular homology $H_*(\mathbb{R}^{2n};\mathbb{Z}_2)$ is concentrated in degree $0$); and (iii) a norm. Standard continuation arguments control the action shift by the oscillation of $\Phi(A) - \Phi(B)$ over the whole space, because the energy estimate involves the entire Floer cylinder; restricting the norm to a compact set $K$ containing the relevant periodic orbits would need an additional argument that we do not have.
\end{remark}

% =========================================================================
% SECTION 3
% =========================================================================
\section{Microlocal Sheaves and Stratified Rank Characteristic Cycles}
\label{sec:microlocal}

We now perform the structural promotion of the realization map to smooth matrix-valued fields $\Phi(A) \in C^\infty(\Omega, \mathrm{Sym}_n(\mathbb{R}))$.

\begin{definition}[\textbf{Rank Stratification}]
\label{def:rank_strat}
Let $\Phi(A) \in C^\infty(\Omega, \mathrm{Sym}_n(\mathbb{R}))$. For each integer $r \in \{0, 1, \dots, n\}$, the rank-$r$ stratum of $\Omega$ is defined by:
\begin{equation}
\Sigma_r(A) \coloneqq \left\{ x \in \Omega \mid \operatorname{rank}(\Phi(A)(x)) = r \right\}.
\end{equation}
\end{definition}

\begin{proposition}[\textbf{OBL-004: Whitney Regularity of Rank Strata}]
\label{prop:whitney_rank_strat}
Let $\Omega$ be a smooth manifold. Assume that the mapping $\Phi(A) \colon \Omega \to \mathrm{Sym}_n(\mathbb{R})$ is transverse to each stratum $\mathcal{R}_r$ of the rank stratification $\mathcal{R} = \{\mathcal{R}_r\}_{r=0}^n$ of $\mathrm{Sym}_n(\mathbb{R})$, where $\mathcal{R}_r$ is the set of symmetric matrices of rank exactly $r$. Then each non-empty $\Sigma_r(A)$ is a submanifold of codimension $\binom{n-r+1}{2}$, and $\Omega = \bigsqcup_{r=0}^n \Sigma_r(A)$ is a Whitney stratification, satisfying conditions~(A) and~(B) of \cite{Whitney1965} (see also \cite[\S\S1--2]{Mather2012}).
\end{proposition}

\begin{proof}
The congruence action $X \mapsto gXg^{\top}$ of $GL_n(\mathbb{R})$ on $\mathrm{Sym}_n(\mathbb{R})$ has finitely many orbits $O_{p,q}$, indexed by the signature $(p,q)$, and $\mathcal{R}_r = \bigsqcup_{p+q=r} O_{p,q}$. Along a continuous path of rank-$r$ matrices no non-zero eigenvalue can cross zero, so each $O_{p,q}$ is open and closed in $\mathcal{R}_r$; hence $\mathcal{R}_r$ is a semi-algebraic submanifold of dimension $r(2n-r+1)/2$, that is, of codimension $\binom{n-r+1}{2}$. For two orbits $O \subset \overline{O'}$, the set of points of $O$ at which Whitney's condition (B) fails relative to $O'$ is semi-algebraic of dimension less than $\dim O$ (Whitney's genericity theorem \cite{Whitney1965}, see \cite{GibsonEtAl1976}) and invariant under the action, which is transitive on $O$; so it is empty. Thus the rank strata form a Whitney stratification of $\mathrm{Sym}_n(\mathbb{R})$. The preimage of a Whitney stratification under a smooth map transverse to all its strata is a Whitney stratification with the same codimensions (see \cite{GibsonEtAl1976}; for the closely related statement on transverse intersections, see \cite{OrroTrotman2010}). Hence $\Sigma_r(A) = \Phi(A)^{-1}(\mathcal{R}_r)$ has codimension $\binom{n-r+1}{2}$ and the partition $\{\Sigma_r(A)\}$ satisfies (A) and (B).
\end{proof}

\begin{lemma}[\textbf{OBL-005: Kashiwara--Schapira Micro-Support Lagrangian Property}]
\label{lem:microsupport_lagrangian}
Assume that $\Omega$ is a real-analytic manifold, that $\Phi(A)$ is real analytic and satisfies the hypothesis of Proposition~\ref{prop:whitney_rank_strat}, so that the strata $\Sigma_r(A)$ are subanalytic. Let $\mathcal{F}_A \in D^b(\Omega)$ be a bounded complex of sheaves of vector spaces whose cohomology sheaves are locally constant of finite rank on each stratum $\Sigma_r(A)$. Then the micro-support $SS(\mathcal{F}_A) \subset T^*\Omega$ is a closed conic subanalytic Lagrangian subset, and
\begin{equation}
SS(\mathcal{F}_A) \subseteq \bigcup_{r=0}^n \overline{T^*_{\Sigma_r(A)} \Omega}.
\end{equation}
\end{lemma}

\begin{proof}
This is a cited result of Kashiwara and Schapira \cite{KashiwaraSchapira1990}. Since the cohomology sheaves of $\mathcal{F}_A$ are locally constant on the strata of a Whitney stratification, the micro-support is contained in the union of the closed conormal bundles of the strata \cite[Prop.~8.4.1]{KashiwaraSchapira1990}. This union is a closed conic isotropic subanalytic set, so $SS(\mathcal{F}_A)$ is isotropic. By the involutivity theorem \cite[Thm.~6.5.4]{KashiwaraSchapira1990}, $SS(\mathcal{F}_A)$ is coisotropic. A closed conic subanalytic set that is both isotropic and coisotropic is Lagrangian; see also \cite[Thm.~8.4.2]{KashiwaraSchapira1990}, as stated in \cite[Thm.~3.7]{Schapira2017}.
\end{proof}

\begin{theorem}[\textbf{OBL-006: Kashiwara's global index formula}]
\label{thm:kashiwara_index_cc}
In the setting of Lemma~\ref{lem:microsupport_lagrangian}, let $\Omega$ be compact without boundary, of dimension $m$. Orient $T^*\Omega$ by $(\dif\theta)^m$, where $\theta$ is the canonical $1$-form, and let $CC(\mathcal{F}_A)$ be the characteristic cycle of $\mathcal{F}_A$ with the orientation conventions of \cite[\S\S2--3]{Kashiwara1985}. Then the Euler characteristic of $\mathcal{F}_A$ is, up to a sign depending only on $m$, the intersection number of its characteristic cycle with the zero section,
\begin{equation}
\chi(\Omega; \mathcal{F}_A) = (-1)^{m(m+1)/2}\, CC(\mathcal{F}_A) \cdot [T^*_\Omega \Omega],
\end{equation}
and, by additivity of the compactly supported Euler characteristic over the finite partition $\{\Sigma_r(A)\}$,
\begin{equation}
\label{eq:chi_strata}
\chi(\Omega; \mathcal{F}_A) = \sum_{r=0}^n \chi_c(\Sigma_r(A))\, \chi\big((\mathcal{F}_A)_{x_r}\big),
\end{equation}
where $x_r$ is any point of $\Sigma_r(A)$ and the sum runs over non-empty strata (if a stratum is disconnected, the term is summed over its connected components, each with its own stalk).
\end{theorem}

\begin{proof}
The first equality is Kashiwara's index theorem \cite[Thm.~4.2]{Kashiwara1985} applied with the constant function $\varphi = 0$, whose differential graph $Y_\varphi$ is the zero section: its hypotheses, that $\{x \in \operatorname{supp}\mathcal{F}_A : \varphi(x) \le t\}$ and $Y_\varphi \cap SS(\mathcal{F}_A)$ be compact, hold because $\Omega$ is compact, and $\mathcal{F}_A$ is constructible by the hypotheses of Lemma~\ref{lem:microsupport_lagrangian}. The sign $(-1)^{m(m+1)/2}$ is the one in \cite[Thm.~4.2]{Kashiwara1985} and depends on the orientation conventions; other normalizations of the orientation of $T^*\Omega$ change it (see \cite[Ch.~IX]{KashiwaraSchapira1990} for the theory in book form). As a check, for the constant sheaf $CC = [T^*_\Omega\Omega]$ \cite[Ex.~8.5(i)]{Kashiwara1985}; with base-then-fibre coordinates $(x,\xi)$ one has $(\dif\theta)^m = m!\,(-1)^{m(m+1)/2}\,\dif x_1\wedge\cdots\wedge\dif x_m\wedge\dif\xi_1\wedge\cdots\wedge\dif\xi_m$, and the self-intersection of the zero section for the base-then-fibre orientation is the Euler number $\chi(\Omega)$ of $T^*\Omega \cong T\Omega$, so the right-hand side equals $\chi(\Omega)$. For \eqref{eq:chi_strata}: for an open $U \subseteq \Omega$ with closed complement $Z$, the exact sequence $\cdots \to H^k_c(U; \mathcal{F}) \to H^k_c(\Omega; \mathcal{F}) \to H^k_c(Z; \mathcal{F}) \to \cdots$ gives $\chi_c(\Omega;\mathcal{F}) = \chi_c(U;\mathcal{F}) + \chi_c(Z;\mathcal{F})$. Since the rank is lower semicontinuous, $\Sigma_r(A)$ is closed in $\Omega \setminus \bigcup_{r' < r} \Sigma_{r'}(A)$; removing the strata one at a time in order of increasing $r$ gives $\chi_c(\Omega;\mathcal{F}_A) = \sum_r \chi_c(\Sigma_r(A); \mathcal{F}_A|_{\Sigma_r(A)})$. On a connected stratum $S$ the cohomology sheaves are locally constant, and $\chi_c(S; \mathcal{L}) = \chi_c(S)\operatorname{rank}\mathcal{L}$ for a local system $\mathcal{L}$ of finite rank on the subanalytic set $S$ (compute with a finite triangulation adapted to $S$). Summing over the cohomology sheaves with signs, and using $\chi = \chi_c$ on compact $\Omega$, gives \eqref{eq:chi_strata}. In the language of constructible functions, \eqref{eq:chi_strata} is the integral of $x \mapsto \chi((\mathcal{F}_A)_x)$ with respect to the Euler characteristic \cite{Viro1988,Schapira1991}.
\end{proof}

\begin{remark}[A formula that fails]
\label{rem:chi_counterexample}
The formula $\chi(\Omega; \mathcal{F}) = \sum_r m_r \chi(\Sigma_r)$, with $m_r$ the coefficients of $CC(\mathcal{F}) = \sum_r m_r [\overline{T^*_{\Sigma_r}\Omega}]$ and the ordinary Euler characteristic $\chi$ of the strata, is false in general. Take $\Omega = S^1$, $n = 1$, $\Phi(A)(\theta) = \cos\theta$ (transverse to $\{0\}$) and $\mathcal{F}$ the constant sheaf. Then $\Sigma_1$ is the union of two open arcs and $\Sigma_0$ consists of two points; $CC(\mathcal{F}) = [T^*_{S^1}S^1] = [\overline{T^*_{\Sigma_1} S^1}]$, so $m_1 = 1$, $m_0 = 0$. The formula would give $1\cdot\chi(\Sigma_1) = 2$, while $\chi(S^1) = 0$. Formula \eqref{eq:chi_strata} gives $\chi_c(\Sigma_1) + \chi_c(\Sigma_0) = -2 + 2 = 0$. In general, the intersection of the closure of a conormal bundle with the zero section involves local Euler obstructions of the stratum closures \cite{MacPherson1974}, not the Euler characteristics of the strata.
\end{remark}

% =========================================================================
% SECTION 4
% =========================================================================
\section{Nonlinear Spectral and Metric Geometry: Weighted \texorpdfstring{$p$}{p}-Laplacians}
\label{sec:nonlinear}

Let $\Omega \subset \mathbb{R}^d$ be a bounded, \emph{connected} domain with Lipschitz boundary $\partial\Omega$, and let $d\mu_A = \Phi(A)(x) \dif x$ with $\Phi(A) \in C^1(\bar{\Omega})$ strictly positive: $\inf_\Omega \Phi(A) \ge c_0 > 0$. For $p \in (1, \infty)$, the weighted $p$-Laplacian operator is defined weakly by
\begin{equation}
\Delta_p^A u \coloneqq \nabla \cdot \left( \Phi(A) \|\nabla u\|^{p-2} \nabla u \right).
\end{equation}

\begin{definition}[\textbf{First $p$-Laplacian Eigenvalue}]
\label{def:p_laplacian}
The first Dirichlet eigenvalue $\lambda_1^{(p)}(A)$ is defined by the Rayleigh--Finsler quotient:
\begin{equation}
\label{eq:rayleigh_p}
\lambda_1^{(p)}(A) \coloneqq \inf_{u \in W_0^{1,p}(\Omega, \mu_A) \setminus \{0\}} \frac{\int_\Omega \|\nabla u\|^p \Phi(A) \dif x}{\int_\Omega |u|^p \Phi(A) \dif x}.
\end{equation}
\end{definition}

\begin{lemma}[\textbf{OBL-007: Existence of a non-negative ground state}]
\label{lem:p_laplacian_minimax}
For every $p \in (1, \infty)$, $\lambda_1^{(p)}(A) > 0$, and the infimum in \eqref{eq:rayleigh_p} is attained by a non-negative function $u_p \in W_0^{1,p}(\Omega)\setminus\{0\}$.
\end{lemma}

\begin{proof}
Write $C_0 \coloneqq \sup_\Omega \Phi(A) < \infty$. Since $c_0 \le \Phi(A) \le C_0$, the weighted quotient lies between $c_0/C_0$ and $C_0/c_0$ times the unweighted quotient $\int\|\nabla u\|^p/\int|u|^p$, and $W_0^{1,p}(\Omega,\mu_A) = W_0^{1,p}(\Omega)$ with equivalent norms. By Poincaré's inequality on the bounded domain $\Omega$, the unweighted quotient is bounded below by a positive constant, so $\lambda_1^{(p)}(A) > 0$. Take a minimizing sequence normalized by $\int|u_k|^p\dif\mu_A = 1$; it is bounded in $W_0^{1,p}(\Omega)$, so a subsequence converges weakly in $W_0^{1,p}$ and, by the Rellich--Kondrachov theorem, strongly in $L^p$. The constraint passes to the limit, and $u \mapsto \int\|\nabla u\|^p\dif\mu_A$ is convex and continuous, hence weakly lower semicontinuous, so the limit $u$ is a minimizer. Since $\|\nabla|u|\| = \|\nabla u\|$ almost everywhere, $u_p \coloneqq |u|$ is also a minimizer.
\end{proof}

\begin{remark}
For $\Phi(A) \equiv 1$, Lindqvist \cite{Lindqvist1990} proved that on any bounded domain the first eigenfunction has no zeros in $\Omega$ and is unique up to multiplication by a constant; see also \cite[Theorem~6]{Lindqvist2008}. We do not prove the weighted analogue here. Strict convexity of $v \mapsto \|\nabla v\|^p$ alone does not give uniqueness, since the Rayleigh quotient is not convex.
\end{remark}

\begin{theorem}[\textbf{OBL-008: The limit $p\to\infty$ is the reciprocal inradius}]
\label{thm:cheeger_p_limit}
Let $R_\Omega \coloneqq \max_{x \in \bar\Omega} \operatorname{dist}(x, \partial\Omega)$ be the inradius of $\Omega$. Then
\begin{equation}
\lim_{p \to \infty} \left( \lambda_1^{(p)}(A) \right)^{1/p} = \frac{1}{R_\Omega},
\end{equation}
for every weight $\Phi(A) \in C^1(\bar\Omega)$ with $\inf_\Omega \Phi(A) > 0$. In particular the limit does not depend on $A$.
\end{theorem}

\begin{proof}
\emph{Step 1: reduction to the unweighted case.} Let $\lambda_p$ denote the unweighted first eigenvalue ($\Phi \equiv 1$). As in the proof of Lemma~\ref{lem:p_laplacian_minimax}, $(c_0/C_0)\lambda_p \le \lambda_1^{(p)}(A) \le (C_0/c_0)\lambda_p$. Since $(C_0/c_0)^{1/p} \to 1$, it suffices to show $\lambda_p^{1/p} \to 1/R_\Omega$. This unweighted statement is due to Juutinen, Lindqvist and Manfredi \cite{JuutinenLindqvistManfredi1999} and to Fukagai, Ito and Narukawa \cite{FukagaiItoNarukawa1999}; we include the short proof for completeness, which follows \cite[Lemma~11]{Lindqvist2008}. Write $\delta(x) = \operatorname{dist}(x, \partial\Omega)$.

\emph{Step 2: upper bound.} The function $\delta$ is $1$-Lipschitz, vanishes on $\partial\Omega$ and satisfies $\|\nabla\delta\| = 1$ almost everywhere, so $\delta \in W_0^{1,p}(\Omega)$ and $\lambda_p^{1/p} \le \|\nabla\delta\|_{L^p}/\|\delta\|_{L^p} = |\Omega|^{1/p}/\|\delta\|_{L^p}$. As $p \to \infty$, $|\Omega|^{1/p} \to 1$ and $\|\delta\|_{L^p} \to \|\delta\|_{L^\infty} = R_\Omega$, so $\limsup_{p} \lambda_p^{1/p} \le 1/R_\Omega$.

\emph{Step 3: lower bound.} Let $L \coloneqq \liminf_{p\to\infty}\lambda_p^{1/p}$, which is finite by Step~2, and choose $p_j \to \infty$ with $\lambda_{p_j}^{1/p_j} \to L$. Let $u_j \ge 0$ be minimizers from Lemma~\ref{lem:p_laplacian_minimax} with $\|u_j\|_{L^{p_j}} = 1$, extended by zero outside $\Omega$. For fixed $q < p_j$, Hölder's inequality gives $\|\nabla u_j\|_{L^q} \le |\Omega|^{1/q - 1/p_j}\lambda_{p_j}^{1/p_j}$ and $\|u_j\|_{L^q} \le |\Omega|^{1/q-1/p_j}$. Fixing $q > d$, the sequence is bounded in $W^{1,q}(\mathbb{R}^d)$ with supports in $\bar\Omega$; by Morrey's inequality and the Arzelà--Ascoli theorem, and a diagonal argument over $q = d+1, d+2, \dots$, a subsequence converges uniformly on $\mathbb{R}^d$ and weakly in every $W^{1,q}(\mathbb{R}^d)$ to a continuous $u_\infty$ vanishing outside $\Omega$. Weak lower semicontinuity gives $\|\nabla u_\infty\|_{L^q} \le |\Omega|^{1/q} L$, and letting $q \to \infty$, $\|\nabla u_\infty\|_{L^\infty} \le L$; hence $u_\infty$ is $L$-Lipschitz on $\mathbb{R}^d$. Moreover $\|u_j\|_{L^\infty} \ge |\Omega|^{-1/p_j}\|u_j\|_{L^{p_j}} = |\Omega|^{-1/p_j} \to 1$, so uniform convergence gives $\|u_\infty\|_{L^\infty} \ge 1$. For $x \in \Omega$ pick $y \in \partial\Omega$ with $|x - y| = \delta(x)$; since $u_\infty(y) = 0$, $|u_\infty(x)| \le L\,\delta(x) \le L R_\Omega$. Taking the supremum over $x$ gives $1 \le L R_\Omega$, that is, $L \ge 1/R_\Omega$.
\end{proof}

\begin{remark}[The Cheeger constant is the limit $p \to 1$, not $p \to \infty$]
\label{rem:cheeger}
For the unit disk and $\Phi(A) \equiv 1$, Theorem~\ref{thm:cheeger_p_limit} gives the limit $1$, whereas the Cheeger constant \cite{Cheeger1970} $h = \inf_E \mathrm{Per}(E)/|E|$ (with the full perimeter) equals $2$, attained by the disk itself. The cone $u = 1 - |x|$ shows directly that $\lambda_p^{1/p} \le \big((p+1)(p+2)/2\big)^{1/p} \to 1$. With the relative perimeter $\int_{\partial^*E \cap \Omega}\Phi(A)\dif\mathcal{H}^{d-1}$, the choice $E = \Omega$ gives a Cheeger quotient $0$. The Cheeger constant appears in the opposite limit: Kawohl and Fridman \cite{KawohlFridman2003} proved, for the unweighted $p$-Laplacian, that $\lambda_p(\Omega) \to h(\Omega)$ as $p \to 1$.
\end{remark}

\begin{theorem}[\textbf{OBL-009: Gromov hyperbolicity of negatively curved conformal metrics}]
\label{thm:gromov_hyperbolicity}
Let $(\Omega, g_A)$ be a complete, simply connected Riemannian manifold with conformal metric $g_A = \Phi(A)^{2/d} \delta_{ij}$. If the sectional curvature satisfies $\mathrm{Sec}(g_A) \le -\kappa_0 < 0$ uniformly on $\Omega$, then $(\Omega, d_{g_A})$ is complete, geodesic, and Gromov $\delta$-hyperbolic with constant:
\begin{equation}
\delta(A) \le \frac{\ln(1 + \sqrt{2})}{\sqrt{\kappa_0}} < \infty.
\end{equation}
\end{theorem}

\begin{proof}
This is a direct consequence of classical comparison geometry. Since $(\Omega, g_A)$ is complete, simply connected and has sectional curvature at most $-\kappa_0 < 0$, the Cartan--Hadamard theorem shows that it is a Hadamard manifold, with unique minimizing geodesics between any two points, and comparison geometry shows that $(\Omega, d_{g_A})$ is a $\mathrm{CAT}(-\kappa_0)$ space \cite{BridsonHaefliger1999}: every geodesic triangle is thinner than its comparison triangle in the model plane $\mathbb{H}^2(-\kappa_0)$. Every geodesic triangle in $\mathbb{H}^2(-\kappa_0)$ is $\delta$-slim with $\delta = \kappa_0^{-1/2}\ln(1+\sqrt2)$, the slimness constant of an ideal triangle (the largest distance from a point of one side to the union of the other two). The comparison map transfers $\delta$-slimness to $(\Omega, d_{g_A})$ \cite{BridsonHaefliger1999,Gromov1987}.
\end{proof}

\begin{remark}
The conformal form $g_A = \Phi(A)^{2/d}\delta_{ij}$ plays no role in the proof; only the curvature bound is used. The constant $\ln(1+\sqrt2) \approx 0.8814$ is not the inradius of the ideal triangle, which is $\tfrac12\ln 3 \approx 0.5493$.
\end{remark}

% =========================================================================
% SECTION 5
% =========================================================================
\section{Non-Equilibrium Thermodynamics and Stochastic Dissipation Length}
\label{sec:thermo}

\emph{Units.} We set $k_B T = 1$ and take unit mobility, so the diffusion coefficient is $D = 1$. Then potentials, work and free energies are dimensionless, and $[x]^2 = [t]$: lengths squared have the dimension of time.

Let $\Phi(A) \in C^2(\mathbb{R}^d, \mathbb{R}_{>0})$ be an unnormalized realization with total mass
\begin{equation}
Z(A) \coloneqq \int_{\mathbb{R}^d} \Phi(A)(x) \dif x = \Tr(A) > 0.
\end{equation}
Define the potential $V_A(x) \coloneqq -\log \Phi(A)(x)$ and consider the overdamped Langevin diffusion on $\mathbb{R}^d$
\begin{equation}
\label{eq:langevin_sde}
\dif X_t = -\nabla V_A(X_t) \dif t + \sqrt{2} \dif W_t.
\end{equation}

\begin{proposition}[\textbf{OBL-010: Stationary measure and $\mathcal{W}_2$ contraction}]
\label{prop:langevin_ergodicity}
Assume that $V_A$ is uniformly convex on $\mathbb{R}^d$: $\nabla^2 V_A(x) \ge \kappa I$ with $\kappa > 0$ for all $x$. Then \eqref{eq:langevin_sde} has a unique strong solution for every initial law with finite second moment, the probability measure
\begin{equation}
\dif\mu_A(x) \coloneqq \frac{e^{-V_A(x)}}{Z(A)} \dif x = \frac{\Phi(A)(x)}{\Tr(A)} \dif x
\end{equation}
is stationary and satisfies the Bakry--\'Emery condition $CD(\kappa, \infty)$, and the law $\mu_t$ of $X_t$ satisfies
\begin{equation}
\mathcal{W}_2(\mu_t, \mu_A) \le e^{-\kappa t}\, \mathcal{W}_2(\mu_0, \mu_A).
\end{equation}
In particular $\mu_A$ is the unique stationary law with finite second moment.
\end{proposition}

\begin{proof}
Uniform convexity gives $\langle \nabla V_A(x) - \nabla V_A(y), x - y\rangle \ge \kappa |x - y|^2$ and $V_A(x) \ge V_A(x_*) + \tfrac{\kappa}{2}|x - x_*|^2$ at the minimizer $x_*$, so $e^{-V_A}$ is integrable with Gaussian tails. The drift is locally Lipschitz and satisfies $\langle -\nabla V_A(x), x - x_*\rangle \le -\kappa|x - x_*|^2$, so strong solutions exist, are unique and do not explode. The density $e^{-V_A}$ solves the stationary Fokker--Planck equation $\nabla\cdot(\nabla \rho + \rho\nabla V_A) = 0$, and the generator $L_A f = \Delta f - \nabla V_A \cdot \nabla f$ is symmetric in $L^2(\mu_A)$; $\Gamma_2(f,f) = \|\nabla^2 f\|^2 + \langle \nabla^2 V_A \nabla f, \nabla f\rangle \ge \kappa\,\Gamma(f,f)$ is the condition $CD(\kappa,\infty)$ \cite{BakryEmery1985}.

For the contraction, let $X_t, Y_t$ solve \eqref{eq:langevin_sde} with the same Brownian motion, with $(X_0, Y_0)$ an optimal coupling of $(\mu_0, \mu_A)$; then $Y_t \sim \mu_A$ for all $t$. The noise cancels in the difference, and
\begin{equation}
\frac{\dif}{\dif t}|X_t - Y_t|^2 = -2\langle \nabla V_A(X_t) - \nabla V_A(Y_t), X_t - Y_t\rangle \le -2\kappa |X_t - Y_t|^2,
\end{equation}
so $|X_t - Y_t| \le e^{-\kappa t}|X_0 - Y_0|$ almost surely. Taking second moments, $\mathcal{W}_2(\mu_t, \mu_A)^2 \le \mathbb{E}|X_t - Y_t|^2 \le e^{-2\kappa t}\mathcal{W}_2(\mu_0,\mu_A)^2$. If $\nu$ is another stationary law with finite second moment, applying this with $\mu_0 = \nu$ gives $\mathcal{W}_2(\nu,\mu_A) \le e^{-\kappa t}\mathcal{W}_2(\nu,\mu_A)$ for all $t$, so $\nu = \mu_A$.
\end{proof}

\noindent For the heat flow on a Riemannian manifold, $\mathcal{W}_2$ contraction estimates characterize lower Ricci bounds, by the theorem of von~Renesse and Sturm~\cite{vonRenesseSturm2005}. The Fokker--Planck equation of \eqref{eq:langevin_sde} is the $\mathcal{W}_2$-gradient flow of the relative entropy $H(\cdot\,|\,\mu_A)$ \cite{JordanKinderlehrerOtto1998}.

\begin{remark}
By the Bakry--\'Emery criterion and the Otto--Villani theorem \cite{BakryEmery1985,OttoVillani2000}, $\mu_A$ also satisfies a logarithmic Sobolev inequality and the Talagrand inequality \cite{Talagrand1996} $\mathcal{W}_2^2(\nu,\mu_A) \le \frac{2}{\kappa}H(\nu\,|\,\mu_A)$; see also \cite[Ch.~22]{Villani2009}. These functional inequalities do not by themselves give the $\mathcal{W}_2$ contraction above, which comes from the coupling argument.
\end{remark}

\begin{theorem}[\textbf{OBL-011: Exact Jarzynski Free Energy Identity on Functional Protocols}]
\label{thm:jarzynski_identity}
Let $s \in [0, 1] \mapsto A_s \in \mathcal{S}_+^n$ be a smooth matrix protocol connecting $A_0 = A$ to $A_1 = B$ over duration $\tau$, with $s(t) = t/\tau$, such that each $V_{A_s}$ satisfies the hypotheses of Proposition~\ref{prop:langevin_ergodicity} and $\int\Phi(A_s)\dif x = \Tr(A_s)$ at $s = 0, 1$. Let $X_t$ solve \eqref{eq:langevin_sde} with the time-dependent potential $V_{A_{s(t)}}$ and $X_0 \sim \mu_A$. The stochastic work $W[X_\cdot] = \int_0^\tau \partial_s V_{A_{s(t)}}(X_t) \dot{s}(t) \dif t$ satisfies the exact Jarzynski equality:
\begin{equation}
\mathbb{E}_{\mu_A} \left[ e^{-W} \right] = e^{-\Delta F(A, B)} = \frac{Z(B)}{Z(A)} = \frac{\Tr(B)}{\Tr(A)}.
\end{equation}
\end{theorem}

\begin{proof}
This is Jarzynski's equality \cite{Jarzynski1997}, applied to the realization potentials; we recall the argument through Crooks' fluctuation theorem \cite{Crooks1999}. Let $P_F[X_\cdot]$ be the forward path measure of the time-dependent diffusion starting from $\mu_A$, and $P_R[\hat{X}_\cdot]$ the path measure of the time-reversed protocol starting from $\mu_B$. By Girsanov's theorem and Crooks' theorem, the Radon--Nikodym derivative satisfies $\ln \frac{\dif P_F}{\dif P_R}[X_\cdot] = W[X_\cdot] - \Delta F(A, B)$, where $\Delta F(A, B) = -\ln(Z(B)/Z(A))$. Integrating the identity $P_R[\hat{X}_\cdot] = e^{-(W - \Delta F)} P_F[X_\cdot]$ over the entire path space yields
\begin{equation}
\mathbb{E}_{\mu_A} \left[ e^{-W} \right] = e^{-\Delta F(A, B)} = \frac{Z(B)}{Z(A)} = \frac{\Tr(B)}{\Tr(A)}.
\end{equation}
\end{proof}

\begin{remark}[\textbf{OBL-012: slow driving, thermodynamic length and $\mathcal{W}_2$}]
\label{rem:thermo_length_w2}
Let $\Sigma_{\mathrm{irr}} \coloneqq \langle W \rangle - \Delta F(A, B) \ge 0$ be the mean dissipated work. In linear response (Sivak and Crooks \cite{SivakCrooks2012}), for a fixed smooth path $s \mapsto A_s$ run at speed $1/\tau$,
\begin{equation}
\tau\,\Sigma_{\mathrm{irr}} \to \int_0^1 g_{\mathrm{fric}}(A_s)(\dot A_s, \dot A_s)\dif s \qquad (\tau\to\infty),
\end{equation}
where $g_{\mathrm{fric}}(A)(\dot A,\dot A) = \int_0^\infty \operatorname{Cov}_{\mu_A}\big(\partial_s V_{A}(X_0), \partial_s V_{A}(X_t)\big)\dif t$ is the friction metric. By the Cauchy--Schwarz inequality the right-hand side is at least $\mathcal{L}^2$, with $\mathcal{L} = \int_0^1 \sqrt{g_{\mathrm{fric}}(\dot A_s,\dot A_s)}\dif s$ the thermodynamic length \cite{SalamonBerry1983,Crooks2007}, and equality holds for constant-speed paths. Thus the slow-driving limit equals $\mathcal{L}^2$, not $\tfrac12\mathcal{L}^2$. We do not prove these linear-response statements here.

A lower bound by $\kappa\,\mathcal{W}_2^2$ fails, and is dimensionally inconsistent: with $k_BT = D = 1$, the quantities $\tau\Sigma_{\mathrm{irr}}$, $\mathcal{L}^2$ and $\mathcal{W}_2^2$ have the dimension of time (length squared), while $\kappa$ has the dimension of inverse time. Counterexample: the dragged trap (the moving laser trap of Schmiedl and Seifert \cite{SchmiedlSeifert2007}, who computed its optimal protocols) $V_s(x) = \kappa(x - sa)^2/2$ on $\mathbb{R}$, $s = t/\tau$, for which $\mu_A$ and $\mu_B$ are Gaussians of equal variance and $\mathcal{W}_2(\mu_A,\mu_B) = |a|$. The mean position obeys $\dot m = -\kappa(m - a t/\tau)$, which gives $\Delta F = 0$ and
\begin{equation}
\tau\,\Sigma_{\mathrm{irr}} = a^2\Big(1 - \frac{1 - e^{-\kappa\tau}}{\kappa\tau}\Big) \longrightarrow a^2 \qquad (\tau\to\infty),
\end{equation}
independently of $\kappa$. Here $g_{\mathrm{fric}}(\dot s) = a^2\dot s^2$, so $\mathcal{L}^2 = a^2$ for this constant-speed protocol; for $\kappa = 10$ the inequality $\lim\tau\Sigma_{\mathrm{irr}} \ge \frac{\kappa}{4}\mathcal{W}_2^2$ would require $a^2 \ge 2.5\,a^2$.

The dimensionally consistent statement is the optimal-transport bound of Aurell, Mej\'ia-Monasterio and Muratore-Ginanneschi \cite{Aurell2011}: over protocols of duration $\tau$ that carry the state $\mu_A$ to the state $\mu_B$, the minimal mean dissipated work is $k_BT\,\mathcal{W}_2^2(\mu_A,\mu_B)/(D\tau)$. In the example above, $\tau\Sigma_{\mathrm{irr}} \to a^2 = \mathcal{W}_2^2$, which matches this minimal value as $\tau \to \infty$. For the relation between thermodynamic length and $\mathcal{W}_2$ in overdamped dynamics, see \cite{NakazatoIto2021,VanVuSaito2023,ZhongDeWeese2024}.
\end{remark}

% =========================================================================
% SECTION 6
% =========================================================================
\section{Bipartite Kernel Operators and Reflected Entropy}
\label{sec:bipartite}

We now realize a positive semi-definite matrix $A \in \mathcal{S}_+^n$ on a bipartite system whose base is the Cartesian product $\Omega \coloneqq \Omega_1 \times \Omega_2$ of two Riemannian manifolds $\Omega_1$ and $\Omega_2$ with their volume measures. The continuous state space is the tensor product Hilbert space:
\begin{equation}
\mathcal{H} \coloneqq L^2(\Omega_1 \times \Omega_2) \cong L^2(\Omega_1) \otimes L^2(\Omega_2).
\end{equation}
Let $\{\psi_k\}_{k=1}^n$ be a fixed orthonormal family in $L^2(\Omega)$ (e.g. Laplacian eigenfunctions or localized wavelet packets). For any matrix $A = (A_{jk}) \in \mathcal{S}_+^n$, its continuous realization is defined algebraically by:
\begin{equation}
\label{eq:kernel_alg_def}
K_A(x, y) \coloneqq \sum_{j, k=1}^n A_{jk} \psi_j(x) \overline{\psi_k(y)}, \qquad x, y \in \Omega_1 \times \Omega_2.
\end{equation}

\begin{proposition}[\textbf{OBL-013: Positive Kernel Codomain Realization and Constructive Diagonal Recovery}]
\label{prop:kernel_codomain_psd}
Let $A \in \mathcal{S}_+^n$ have unitary spectral decomposition $A = U \operatorname{diag}(\lambda_1, \dots, \lambda_n) U^\dagger$. Define the rotated functional basis:
\begin{equation}
\phi_j(x) \coloneqq \sum_{k=1}^n U_{kj} \psi_k(x), \qquad j = 1, \dots, n.
\end{equation}
Then $\{\phi_j\}_{j=1}^n$ is an orthonormal family in $L^2(\Omega_1 \times \Omega_2)$, the kernel $K_A$ in \eqref{eq:kernel_alg_def} is identically diagonalized:
\begin{equation}
K_A(x, y) = \sum_{j=1}^n \lambda_j \phi_j(x) \overline{\phi_j(y)},
\end{equation}
and the induced integral operator $T_A$ is self-adjoint, trace-class, positive semi-definite with $\Tr(T_A) = \Tr(A)$, whose diagonal restriction constructively coincides with the scalar realization $\Phi_{\mathrm{scalar}}(A)(x) \coloneqq \sum_{j,k} A_{jk} \psi_j(x) \overline{\psi_k(x)}$.
\end{proposition}

\begin{proof}
Orthonormality of $\{\phi_j\}$ is guaranteed by the unitarity of $U$:
\begin{equation}
\langle \phi_i, \phi_j \rangle_{L^2} = \sum_{k, l=1}^n U_{ki}^* U_{lj} \langle \psi_k, \psi_l \rangle_{L^2} = \sum_{k=1}^n U_{ki}^* U_{kj} = (U^\dagger U)_{ij} = \delta_{ij}.
\end{equation}
Expanding $K_A(x, y)$ in terms of $\lambda_j$ and $\phi_j$, we compute $\langle f, T_A f \rangle = \sum \lambda_j |\langle f, \phi_j \rangle|^2 \ge 0$ since $\lambda_j \ge 0$. The trace is $\int_\Omega K_A(x, x) \dif x = \sum \lambda_j \|\phi_j\|^2 = \sum \lambda_j = \Tr(A)$. Finally, $K_A(x, x) = \sum \lambda_j |\phi_j(x)|^2 = \sum_{j,k} A_{jk} \psi_j(x) \overline{\psi_k(x)} = \Phi_{\mathrm{scalar}}(A)(x)$ holds identically by matrix multiplication.
\end{proof}

\begin{proposition}[\textbf{OBL-014: Partial trace and modular Hamiltonian}]
\label{thm:modular_hamiltonian_spec}
On $\Omega = \Omega_1 \times \Omega_2$, let $\rho_{\Omega_1}$ be the operator on $L^2(\Omega_1)$ with kernel
\begin{equation}
\rho_{\Omega_1}(x_1, y_1) \coloneqq \int_{\Omega_2} K_A(x_1, z_2; y_1, z_2) \dif z_2 .
\end{equation}
Assume $A \ne 0$.
\begin{enumerate}
\item $\rho_{\Omega_1}$ is self-adjoint, positive semi-definite and trace class, with $\Tr\rho_{\Omega_1} = \Tr(A)$.
\item If each $\psi_k$ has finite Schmidt rank \cite{Schmidt1907} $m_k$, that is, $\psi_k = \sum_{l=1}^{m_k} a_{kl}\otimes b_{kl}$ with $a_{kl} \in L^2(\Omega_1)$, $b_{kl} \in L^2(\Omega_2)$, then $\operatorname{rank}\rho_{\Omega_1} \le \sum_{k=1}^n m_k$. In particular, for a product frame $\psi_k = a_k \otimes b_k$, $\operatorname{rank}\rho_{\Omega_1} \le n$.
\item Let $\hat\rho \coloneqq \rho_{\Omega_1}/\Tr(A)$, with non-zero eigenvalues $p_1 \ge p_2 \ge \dots > 0$ (finitely or countably many, $\sum_k p_k = 1$) and orthonormal eigenvectors $e_k$ spanning $\mathcal{H}_{\mathrm{supp}}^{(1)} \coloneqq (\ker \rho_{\Omega_1})^\perp$. The modular Hamiltonian
\begin{equation}
K_{\Omega_1} \coloneqq -\log\big(\hat\rho|_{\mathcal{H}_{\mathrm{supp}}^{(1)}}\big) = \sum_k (-\log p_k)\,|e_k\rangle\langle e_k|
\end{equation}
is self-adjoint and non-negative on $\mathcal{H}_{\mathrm{supp}}^{(1)}$, with spectrum $\{E_k = -\log p_k\}$. It is strictly positive if and only if $\operatorname{rank}\rho_{\Omega_1} \ge 2$.
\item Let $\Omega_1$, $\Omega_2$ be compact, let $k_{\Omega_i}$ be continuous positive semi-definite kernels with $k_{\Omega_2} \ne 0$, and assume that the integral operator $T_{k_{\Omega_1}}$ is injective on $L^2(\Omega_1)$. For $\epsilon > 0$ and $K_A^\epsilon \coloneqq K_A + \epsilon\,(k_{\Omega_1} \otimes k_{\Omega_2})$, the partial trace $\rho^\epsilon_{\Omega_1}$ is injective. If $L^2(\Omega_1)$ is infinite-dimensional, the modular Hamiltonian $K^\epsilon_{\Omega_1}$ of the normalized $\rho^\epsilon_{\Omega_1}$ is defined on a dense domain of $L^2(\Omega_1)$ and has purely discrete spectrum $E_k^\epsilon \to +\infty$.
\end{enumerate}
\end{proposition}

\begin{proof}
(1) By Proposition~\ref{prop:kernel_codomain_psd}, $T_A = \sum_{j} \lambda_j |\phi_j\rangle\langle\phi_j|$ with $\lambda_j \ge 0$. For $v \in L^2(\Omega_1\times\Omega_2)$ let $V \colon L^2(\Omega_2) \to L^2(\Omega_1)$ be the Hilbert--Schmidt operator with kernel $v(x_1, z_2)$. The operator with kernel $\int v(x_1,z_2)\overline{v(y_1,z_2)}\dif z_2$ is $VV^*$, which is positive semi-definite and trace class with $\Tr VV^* = \|V\|_{HS}^2 = \|v\|^2$. Hence $\rho_{\Omega_1} = \sum_j \lambda_j \Phi_j\Phi_j^*$, where $\Phi_j$ has kernel $\phi_j(x_1,z_2)$, is positive semi-definite and trace class with $\Tr\rho_{\Omega_1} = \sum_j \lambda_j\|\phi_j\|^2 = \Tr(A)$.

(2) Each $\phi_j = \sum_k U_{kj}\psi_k$ lies in $\operatorname{span}\{a_{kl}\otimes b_{kl}\}$, so the range of $\Phi_j$ lies in $W \coloneqq \operatorname{span}\{a_{kl}\}$, of dimension at most $\sum_k m_k$. The range of $\rho_{\Omega_1} = \sum_j\lambda_j\Phi_j\Phi_j^*$ lies in $W$.

(3) By (1), $\hat\rho$ is a positive trace-class operator of trace $1$, so $0 < p_k \le 1$ and $E_k = -\log p_k \ge 0$. If the rank is $1$, then $p_1 = 1$ and $E_1 = 0$. If the rank is at least $2$, then every $p_k < 1$, so every $E_k > 0$, and $E_k \to +\infty$ if there are infinitely many.

(4) The partial trace of $k_{\Omega_1}\otimes k_{\Omega_2}$ is $c\,T_{k_{\Omega_1}}$ with $c = \int_{\Omega_2} k_{\Omega_2}(z,z)\dif z > 0$, and $T_{k_{\Omega_i}}$ is trace class by Mercer's theorem \cite{Mercer1909}. Hence $\rho^\epsilon_{\Omega_1} = \rho_{\Omega_1} + \epsilon c\,T_{k_{\Omega_1}} \ge \epsilon c\,T_{k_{\Omega_1}}$. If $\rho^\epsilon_{\Omega_1} f = 0$, then $0 = \langle f, \rho^\epsilon_{\Omega_1} f\rangle \ge \epsilon c\,\|T_{k_{\Omega_1}}^{1/2} f\|^2$, so $T_{k_{\Omega_1}} f = 0$ and $f = 0$. A compact, injective, positive self-adjoint operator on an infinite-dimensional space has an orthonormal eigenbasis with eigenvalues $p_k^\epsilon \to 0$; so $-\log$ of its normalization is a self-adjoint operator, diagonal in that basis, with eigenvalues $E_k^\epsilon \to +\infty$, hence with compact resolvent.
\end{proof}

\begin{remark}[Two statements that fail]
\label{rem:rank_counterexample}
(a) Without the finite Schmidt-rank hypothesis in (2), $\rho_{\Omega_1}$ can have infinite rank even for $n = 1$. Take $\Omega_1 = \Omega_2 = \mathbb{R}$ and $\psi_1(x, y) \propto \exp\big(-(x^2+y^2)/2 + \beta x y\big)$ with $0 < |\beta| < 1$. By Mehler's formula \cite[eq.~18.18.28]{DLMF} its Schmidt coefficients are geometric, $p_k \propto q^k$ with $q = \big((1 - \sqrt{1-\beta^2})/\beta\big)^2$, so all are non-zero. For $\beta = 0.6$, $q = 1/9$. The same geometric spectrum appears for the reduced ground state of two coupled harmonic oscillators \cite{Srednicki1993}.
(b) Without normalization, $-\log$ of the restriction of $\rho_{\Omega_1}$ need not be non-negative: for $n = 1$, $A = (5)$ and a product frame, $\rho_{\Omega_1} = 5\,|a_1\rangle\langle a_1|$ and $-\log 5 < 0$. Even after normalization, rank $1$ gives the eigenvalue $0$.
\end{remark}

\begin{theorem}[\textbf{OBL-015: Canonical Purification, Symmetry and the Mutual-Information Bound for the Reflected Entropy}]
\label{thm:reflected_entropy_prop}
Assume that $\rho_{\Omega_1}$ and $\rho_{\Omega_2}$ have finite rank (for instance, a frame of finite Schmidt ranks, Proposition~\ref{thm:modular_hamiltonian_spec}(2)), so that $\mathcal{H}_{\mathrm{supp}}^{(1)}$ and $\mathcal{H}_{\mathrm{supp}}^{(2)}$ are finite-dimensional. Let $\rho_{12} \coloneqq T_A/\Tr(A)$, a density operator supported on $\mathcal{H}_{\mathrm{supp}}^{(1)} \otimes \mathcal{H}_{\mathrm{supp}}^{(2)}$, and let $|\sqrt{\rho_{12}}\rangle \in (\mathcal{H}_{\mathrm{supp}}^{(1)} \otimes \mathcal{H}_{\mathrm{supp}}^{(2)}) \otimes (\mathcal{H}_{\mathrm{supp}}^{(1^*)} \otimes \mathcal{H}_{\mathrm{supp}}^{(2^*)})$ be its canonical GNS purification. The reflected entropy
\begin{equation}
S_R(\Omega_1 : \Omega_2) \coloneqq -\Tr\left( \sigma_{1 1^*} \log \sigma_{1 1^*} \right), \qquad \sigma_{1 1^*} \coloneqq \Tr_{2 2^*}\left( |\sqrt{\rho_{12}}\rangle \langle\sqrt{\rho_{12}}| \right)
\end{equation}
is non-negative, symmetric $S_R(\Omega_1 : \Omega_2) = S_R(\Omega_2 : \Omega_1)$, and satisfies the universal lower bound:
\begin{equation}
S_R(\Omega_1 : \Omega_2) \ge I(\Omega_1 : \Omega_2) \coloneqq S(\Omega_1) + S(\Omega_2) - S(\Omega_1 \Omega_2).
\end{equation}
No monotonicity of $S_R$ under partial trace is asserted: Hayden, Lemm and Sorce \cite[Thm.~1]{HaydenLemmSorce2023} give classical states on $\mathbb{C}^3\otimes\mathbb{C}^3\otimes\mathbb{C}^2$ with $S_R(A:BC) < S_R(A:B)$.
\end{theorem}

\begin{proof}
All spaces are finite-dimensional. $S_R$ is a von Neumann entropy, hence non-negative. The vector $|\sqrt{\rho_{12}}\rangle$ is a pure state on $11^*22^*$, so the reduced states on $11^*$ and on $22^*$ have the same non-zero spectrum, which gives the symmetry. The bound $S_R \ge I(\Omega_1 : \Omega_2)$ is \cite[eqs.~(2.12)--(2.14)]{DuttaFaulkner2019}: since the state on $11^*22^*$ is pure, $S(11^*2) = S(2^*) = S(2)$, so $S_R = S(11^*) = I(11^* : 2)$, and $I(11^*:2) \ge I(1:2)$ by monotonicity of mutual information under discarding $1^*$, which is strong subadditivity \cite{LiebRuskai1973}.
\end{proof}

\begin{remark}[\textbf{OBL-016: the holographic proposal}]
\label{thm:holographic_ew_inequality}
For holographic quantum field theories with a classical gravity dual, Dutta and Faulkner \cite{DuttaFaulkner2019} proposed, and supported by a gravitational replica argument, that at leading order in $G_N$
\begin{equation}
S_R(A : B) = 2 E_W(A : B), \qquad E_W = \frac{\min_{\Sigma_{AB}} \mathrm{Area}(\Sigma_{AB})}{4 G_N},
\end{equation}
where $\Sigma_{AB}$ runs over surfaces splitting the entanglement wedge of $A \cup B$ into a part adjacent to $A$ and a part adjacent to $B$ (in units $\hbar = c = 1$, so that $\mathrm{Area}/G_N$ is dimensionless). The cross-section $E_W$ had been proposed earlier as the holographic dual of the entanglement of purification \cite{UmemotoTakayanagi2018,NguyenEtAl2018}; for the gap between $S_R$ and the mutual information in holographic states see \cite{AkersRath2020,HaydenParrikarSorce2021}. This is a physical duality, not a theorem. No bulk geometry or boundary state has been constructed here from the kernel $K_A$, so the relation is not established for the realizations of this section; we record it only as motivation.
\end{remark}

% =========================================================================
% SECTION 7
% =========================================================================
\section{Mellin Transforms over the Simplex}
\label{sec:simplicial}

Let $\Delta_m \coloneqq \{ u \in \mathbb{R}_{\ge 0}^m \mid \sum_{i=1}^m u_i = 1 \}$ be the standard simplex, equipped with the Dirichlet probability measure $\dif\mu_\alpha(u) = \frac{\Gamma(\alpha_0)}{\prod_i \Gamma(\alpha_i)} \prod_i u_i^{\alpha_i - 1} \dif u$, where $\alpha_i > 0$, $\alpha_0 \coloneqq \sum_i \alpha_i$, and $\dif u$ is the Lebesgue measure on the first $m-1$ coordinates.

\begin{proposition}[\textbf{OBL-017: The simplicial Mellin transform is entire}]
\label{prop:simplicial_mellin_conv}
For any continuous, strictly positive realization $\Phi(A)$ on $\Delta_m$, the simplicial Mellin transform
\begin{equation}
\mathcal{Z}_A(s) \coloneqq \int_{\Delta_m} \left( \Phi(A)(u) \right)^s \dif\mu_\alpha(u)
\end{equation}
is an entire function of $s \in \mathbb{C}$. In particular it has no poles.
\end{proposition}

\begin{proof}
On the compact simplex, $0 < c_1 \le \Phi(A) \le c_2 < \infty$. For each $u$, $s \mapsto \Phi(A)(u)^s = e^{s\log\Phi(A)(u)}$ is entire, and for $|s| \le S$ we have $|\Phi(A)(u)^s| \le e^{S\,|\log\Phi(A)(u)|} \le e^{S \max(|\log c_1|, |\log c_2|)}$, a bound independent of $u$. Since $\mu_\alpha$ is a probability measure, $\mathcal{Z}_A$ is continuous by dominated convergence, and Fubini's and Morera's theorems show that it is holomorphic on every disk $|s| < S$, hence entire.
\end{proof}

\noindent For an affine realization $\Phi(A)(u) = \sum_i z_i u_i$ with all $z_i > 0$, $\mathcal{Z}_A(s)$ is the Dirichlet average of the power $s$, that is, Carlson's $R$-function $R_s(\alpha; z)$ \cite{Carlson1963}, a symmetric form of the Lauricella function $F_D$; see \cite[\S3.3.1]{Chow2022}.

\begin{remark}[\textbf{OBL-018: no residue spectrum for $\mathcal{Z}_A$}]
\label{thm:barnes_meromorphic_continuation}
By Proposition~\ref{prop:simplicial_mellin_conv}, $\mathcal{Z}_A$ has no poles, so it has no residue spectrum at $s = -\alpha_j - k$. The simplest instance is $\Phi(A) \equiv c > 0$, for which $\mathcal{Z}_A(s) = c^s$. Poles at $s = -\alpha_j - k$ belong instead to Mellin transforms of a vertex coordinate, whose integrand is not bounded away from $0$; see Proposition~\ref{prop:coordinate_mellin}.
\end{remark}

\begin{proposition}[Mellin transform of a vertex coordinate]
\label{prop:coordinate_mellin}
For $j \in \{1, \dots, m\}$ and $m \ge 2$, let $M_j(s) \coloneqq \int_{\Delta_m} u_j^s \dif\mu_\alpha(u)$, which converges for $\operatorname{Re}(s) > -\alpha_j$. Then
\begin{equation}
M_j(s) = \frac{\Gamma(\alpha_0)\,\Gamma(\alpha_j + s)}{\Gamma(\alpha_j)\,\Gamma(\alpha_0 + s)},
\end{equation}
which extends meromorphically to $\mathbb{C}$. Write $b_j \coloneqq \alpha_0 - \alpha_j > 0$. If $b_j \notin \mathbb{Z}$, the poles are exactly $s = -\alpha_j - k$, $k \in \mathbb{N}_0$, all simple, with residues
\begin{equation}
\operatorname{Res}_{s = -\alpha_j - k} M_j = \frac{(-1)^k}{k!}\,\frac{\Gamma(\alpha_0)}{\Gamma(\alpha_j)\,\Gamma(b_j - k)} \ne 0.
\end{equation}
If $b_j$ is a positive integer, $M_j$ is a rational function with simple poles exactly at $s = -\alpha_j - k$, $0 \le k \le b_j - 1$.
\end{proposition}

\begin{proof}
Under $\mu_\alpha$ the coordinate $u_j$ has the Beta$(\alpha_j, b_j)$ distribution (aggregation property of the Dirichlet distribution), so for $\operatorname{Re}(s) > -\alpha_j$,
$M_j(s) = B(\alpha_j + s, b_j)/B(\alpha_j, b_j)$, which is the stated formula by the Euler beta integral \cite[eq.~5.12.1]{DLMF}. The function $1/\Gamma$ is entire, so the right-hand side is meromorphic with possible poles only at the poles $s = -\alpha_j - k$ of $\Gamma(\alpha_j + s)$, where $\operatorname{Res}\Gamma(\alpha_j+s) = (-1)^k/k!$ \cite[\S5.2(i)]{DLMF}. At such a point $\Gamma(\alpha_0 + s) = \Gamma(b_j - k)$, which is finite and non-zero unless $b_j - k \in \{0, -1, -2, \dots\}$. If $b_j \notin \mathbb{Z}$ this never happens, which gives the residues. If $b_j \in \mathbb{Z}_{>0}$, then $\Gamma(\alpha_j+s)/\Gamma(\alpha_j+b_j+s) = \prod_{i=0}^{b_j-1}(\alpha_j+s+i)^{-1}$.
\end{proof}

% =========================================================================
% SECTION 8
% =========================================================================
\section{Federer Reach and Medial Axis Geometry of Level Surfaces}
\label{sec:federer}

\begin{theorem}[\textbf{OBL-020: Level Set Federer Reach Lower Bound}]
\label{thm:federer_reach_bound}
Let $\Omega \subseteq \mathbb{R}^d$ be open, $d \ge 2$, and $\Phi(A) \in C^\infty(\Omega, \mathbb{R})$. Let $t$ be a regular value such that $\Sigma_t = \Phi(A)^{-1}(t)$ is non-empty and compact, and set $\epsilon_0 \coloneqq \min_{\Sigma_t} \|\nabla \Phi(A)\| > 0$, $M \coloneqq \max_{\Sigma_t} \|\nabla^2 \Phi(A)\|_{\mathrm{op}}$, and $\nu \coloneqq \nabla\Phi(A)/\|\nabla\Phi(A)\|$ on $\Sigma_t$. Let
\begin{equation}
d_{\mathrm{sep}}(\Sigma_t) \coloneqq \inf \{ \|x - y\| \mid x, y \in \Sigma_t,\ \nu(x) = -\nu(y) \}.
\end{equation}
Then $d_{\mathrm{sep}}(\Sigma_t) > 0$, and the reach of $\Sigma_t$ satisfies (with $\epsilon_0/M = \infty$ if $M = 0$)
\begin{equation}
\mathrm{reach}(\Sigma_t) \ge \min\left( \frac{\epsilon_0}{M}, \frac{1}{2} d_{\mathrm{sep}}(\Sigma_t) \right) > 0.
\end{equation}
\end{theorem}

\begin{proof}
$\Sigma_t$ is a compact embedded smooth hypersurface of $\mathbb{R}^d$. By uniform continuity of $\nu$ there is $\eta > 0$ with $\|\nu(x) - \nu(y)\| < 2$ whenever $\|x-y\| < \eta$, so $d_{\mathrm{sep}} \ge \eta > 0$.

\emph{Curvature.} Let $\gamma$ be a unit-speed geodesic of $\Sigma_t$ with $\gamma(0) = x$, $\gamma'(0) = v$. Differentiating $\langle \nabla\Phi(A)(\gamma), \gamma'\rangle = 0$ gives $\nabla^2\Phi(A)(v, v) + \langle\nabla\Phi(A), \gamma''(0)\rangle = 0$; since $\gamma''(0)$ is normal to $\Sigma_t$, $\|\gamma''(0)\| = |\nabla^2\Phi(A)(v,v)|/\|\nabla\Phi(A)\| \le M/\epsilon_0$.

\emph{Reach.} The reach $\tau$ of a compact smooth submanifold is positive, by Federer's formula $\tau = \inf_{p\ne q}\|q-p\|^2/(2\operatorname{dist}(q-p, T_p\Sigma_t))$ \cite[Thm.~4.18]{Federer1959} and compactness. By the dichotomy of Aamari et al.\ \cite[Thm.~3.4]{Aamari2019}, either (local case) some unit-speed geodesic has $\|\gamma''(0)\| = 1/\tau$, and then $\tau \ge \epsilon_0/M$ by the curvature bound (compare \cite[Prop.~6.1]{NiyogiSmaleWeinberger2008}); or (global case) there are $q_1, q_2 \in \Sigma_t$ and $z_0$ with $q_1, q_2 \in \partial B(z_0, \tau)$, $\|q_1 - q_2\| = 2\tau$, and the open ball $B(z_0, \tau)$ disjoint from $\Sigma_t$. In the global case $z_0$ is the midpoint of $q_1 q_2$, the segment is normal to $\Sigma_t$ at both ends, and $\Phi(A) - t$ has constant sign on the connected ball, say negative. Then $\Phi(A)$ decreases from $q_i$ towards $z_0$, so $\nu(q_i) = (q_i - z_0)/\tau$ for $i = 1, 2$; hence $\nu(q_1) = -\nu(q_2)$ and $2\tau = \|q_1 - q_2\| \ge d_{\mathrm{sep}}$. (If $\Phi(A) - t > 0$ on the ball, both normals change sign and the conclusion is the same.) In both cases $\tau \ge \min(\epsilon_0/M, d_{\mathrm{sep}}/2)$.
\end{proof}

\begin{remark}
\label{rem:reach_examples}
(a) Compactness cannot be dropped: for $\Phi(A)(x,y) = y$ on $\Omega = (0,1)^2$, the level set $\Sigma_t$ is an open segment, with $M = 0$ and no pair of opposite normals, so the right-hand side is $+\infty$, while the reach of the (non-closed) segment is $0$.
(b) For a strictly convex $K_t(A)$ the quantity $d_{\mathrm{sep}}$ is finite, since the Gauss map of the boundary is onto the sphere. For the ellipse $\Phi(A)(x,y) = x^2/4 + y^2$, $t = 1$: $\epsilon_0 = 1$, $M = 2$, opposite normals occur exactly at antipodal pairs $\pm x$, so $d_{\mathrm{sep}} = 2$, and the bound gives $\min(1/2, 1) = 1/2$, which equals the reach (the minimal radius of curvature $b^2/a = 1/2$).
(c) Federer \cite[4.21]{Federer1959} records a lower bound for the reach of a preimage $f^{-1}(A)$ in terms of the Lipschitz constants of $f$ and $Df$ and a lower bound for the Jacobian of $f$ on a neighbourhood of $f^{-1}(A)$. Theorem~\ref{thm:federer_reach_bound} uses only $\Phi(A)$ and its first two derivatives on $\Sigma_t$, together with $d_{\mathrm{sep}}$. The reach of the obstacle sets in Chapter~7 of~\cite{silvafilho2026book} is treated with the same notion.
\end{remark}

\begin{theorem}[\textbf{OBL-021: Medial Axis Avoidance and Weyl--Federer Tube Volume Expansion}]
\label{thm:tube_volume_reach}
In the setting of Theorem~\ref{thm:federer_reach_bound}, let $R_t \coloneqq \mathrm{reach}(\Sigma_t) > 0$ and $U_r(\Sigma_t) \coloneqq \{y \in \mathbb{R}^d : \operatorname{dist}(y, \Sigma_t) < r\}$. Let $\mathcal{M}(\Sigma_t)$ be the medial axis, the set of points of $\mathbb{R}^d$ with at least two nearest points on $\Sigma_t$. For any $r < R_t$, $U_r(\Sigma_t) \cap \mathcal{M}(\Sigma_t) = \emptyset$, the metric projection $\Proj_{\Sigma_t} \colon U_r \to \Sigma_t$ is single-valued and Lipschitz, and
\begin{equation}
\mathrm{Vol}_d(U_r(\Sigma_t)) = \sum_{k=0}^{\lfloor \frac{d-1}{2} \rfloor} \frac{2 r^{2k+1}}{2k+1} \int_{\Sigma_t} H_{2k}(x) \dif\mathcal{H}^{d-1}(x),
\end{equation}
where $H_{j} = \sum_{i_1 < \dots < i_j} \kappa_{i_1}\cdots\kappa_{i_j}$ is the (unnormalized) $j$-th elementary symmetric function of the principal curvatures $\kappa_1, \dots, \kappa_{d-1}$, and $H_0 = 1$.
\end{theorem}

\begin{proof}
By definition of the reach, every point at distance less than $R_t$ from $\Sigma_t$ has a unique nearest point, so $U_r \cap \mathcal{M} = \emptyset$; the projection is Lipschitz on $U_r$ for $r < R_t$ by \cite[Thm.~4.8]{Federer1959}. The map $F(x, s) = x + s\nu(x)$ sends $\Sigma_t \times (-r, r)$ onto $U_r$: a point $y \in U_r$ satisfies $y = F(\Proj y, \pm\operatorname{dist}(y,\Sigma_t))$, because $y - \Proj y$ is normal to $\Sigma_t$ at $\Proj y$. It is injective, because $\Proj(F(x,s)) = x$ for $|s| < R_t$ \cite[Thm.~4.8]{Federer1959}. Its Jacobian is $\prod_{i}(1 - s\kappa_i(x))$, which is positive since the second fundamental form of a closed submanifold of reach $R_t$ satisfies $\|\mathrm{II}(v,v)\| \le 1/R_t$ for unit $v$ \cite[Prop.~A.1]{Aamari2019}, (see also \cite[Prop.~6.1]{NiyogiSmaleWeinberger2008}), so $|\kappa_i| \le 1/R_t$. Hence
$\mathrm{Vol}_d(U_r) = \int_{\Sigma_t}\int_{-r}^{r}\prod_i(1 - s\kappa_i)\dif s\dif\mathcal{H}^{d-1} = \sum_j (-1)^j \int_{\Sigma_t} H_j \int_{-r}^{r}s^j\dif s\dif\mathcal{H}^{d-1}$. The odd terms vanish and $\int_{-r}^r s^{2k}\dif s = 2r^{2k+1}/(2k+1)$.
\end{proof}

\noindent This is the hypersurface case of Weyl's tube formula \cite{Weyl1939}. Federer's Steiner formula \cite[Thm.~5.6]{Federer1959} extends the polynomial dependence on $r$, for $0 \le r < \mathrm{reach}$, to all sets of positive reach.

% =========================================================================
% SECTION 9
% =========================================================================
\section{Summary of Invariants}
\label{sec:taxonomy}

Table~\ref{tab:taxonomy} lists the invariants of this volume, the kind of object each one is, and where it is treated.

\begin{table}[htbp]
\centering
\small
\caption{Invariants of functional realizations treated in this volume.}
\label{tab:taxonomy}
\begin{tabular}{llll}
\toprule
\textbf{Invariant} & \textbf{Type of value} & \textbf{Structure} & \textbf{Where} \\
\midrule
$c_1^{\mathrm{EH}}(K_t(A)), c_{\mathrm{HZ}}(K_t(A))$ & number in $(0,\infty)$ & symplectic capacity & Sec.~\ref{sec:symplectic} \\
$CC(\mathcal{F}_A)$ & Lagrangian cycle in $T^*\Omega$ & microlocal sheaf theory & Sec.~\ref{sec:microlocal} \\
$\lambda_1^{(p)}(A)$ & number in $(0,\infty)$ & weighted $p$-Laplacian & Sec.~\ref{sec:nonlinear} \\
$\delta(A)$ & number in $[0,\infty)$ & Gromov hyperbolicity & Sec.~\ref{sec:nonlinear} \\
$\Delta F(A,B)$, $\Sigma_{\mathrm{irr}}$ & number & stochastic thermodynamics & Sec.~\ref{sec:thermo} \\
$\rho_{\Omega_1}$, $K_{\Omega_1}$, $S_R$ & operator / number & bipartite kernels & Sec.~\ref{sec:bipartite} \\
$\mathcal{Z}_A(s)$ & entire function of $s$ & Mellin transform on $\Delta_m$ & Sec.~\ref{sec:simplicial} \\
$\mathrm{reach}(\Sigma_t)$ & number in $(0,\infty]$ & geometric measure theory & Sec.~\ref{sec:federer} \\
\bottomrule
\end{tabular}
\end{table}

\section*{Declarations}

\noindent\textbf{Funding.} This research was financed in part by the
Coordena\c{c}\~ao de Aperfei\c{c}oamento de Pessoal de N\'ivel Superior -- Brasil
(CAPES) -- Finance Code 001.

\medskip\noindent\textbf{Competing interests.} The author declares no competing
interests.

\medskip\noindent\textbf{Use of generative AI tools.} Large language model
assistants (Anthropic Claude; Google Gemini, through the Antigravity agent
environment) were used extensively. Under the author's direction they drafted and
edited text and \LaTeX{} source, proposed and wrote derivations, proofs and
counterexamples, wrote the numerical check scripts, and searched the literature.
The mathematics was then checked in three steps by AI sessions that had not
written it: a blind review of every statement, a correction made by a separate
session, and a targeted re-check of each correction. Every numerical claim is
checked by a script that compares it with an independent computation and includes
a negative control, a deliberately altered formula that must fail. Every reference
was resolved through Crossref, DataCite or arXiv, or, for the three items without
a DOI, through Numdam, EuDML or the ISBN record. These checks were carried out by
AI systems; they are not peer review, and no result in this volume has yet
been checked by an independent human expert. AI tools are not authors. The author
chose the questions, directed the work, decided what to publish, and is
responsible for the content, including any remaining errors.

\medskip\noindent\textbf{Verification status.} Results labelled Theorem,
Proposition or Lemma are proved in the text or cited as classical results;
statements labelled Conjecture are not proved. The numerical scripts are checks,
not proofs. The \textsc{Lean 4} files in the author's repository associated with
this work are a placeholder skeleton without Mathlib and verify none of the
statements.

\medskip\noindent\textbf{Code availability.} The numerical check scripts are
available from the author on request.

% =========================================================================
% REFERENCES
% =========================================================================
\begin{thebibliography}{99}

\bibitem{Aamari2019}
E.~Aamari, J.~Kim, F.~Chazal, B.~Michel, A.~Rinaldo, and L.~Wasserman,
\emph{Estimating the reach of a manifold},
Electronic Journal of Statistics, \textbf{13} (2019),
\href{https://doi.org/10.1214/19-EJS1551}{DOI: 10.1214/19-EJS1551}.

\bibitem{Aurell2011}
E.~Aurell, C.~Mej\'ia-Monasterio, and P.~Muratore-Ginanneschi,
\emph{Optimal protocols and optimal transport in stochastic thermodynamics},
Physical Review Letters, \textbf{106} (2011), no.~25, 250601,
\href{https://doi.org/10.1103/PhysRevLett.106.250601}{DOI: 10.1103/PhysRevLett.106.250601}.

\bibitem{BakryEmery1985}
D.~Bakry and M.~\'Emery,
\emph{Diffusions hypercontractives},
S\'eminaire de Probabilit\'es XIX 1983/84, Lecture Notes in Mathematics, vol.~1123, Springer, Berlin, Heidelberg, 1985, pp.~177--206,
ISBN: 978-3-540-15230-9,
\href{https://doi.org/10.1007/BFb0075847}{DOI: 10.1007/BFb0075847}.

\bibitem{Crooks1999}
G.~E. Crooks,
\emph{Entropy production fluctuation theorem and the nonequilibrium work relation for free energy differences},
Physical Review E, \textbf{60} (1999), no.~3, 2721--2726,
\href{https://doi.org/10.1103/PhysRevE.60.2721}{DOI: 10.1103/PhysRevE.60.2721}.

\bibitem{DuttaFaulkner2019}
S.~Dutta and T.~Faulkner,
\emph{A canonical purification for the entanglement wedge cross-section},
Journal of High Energy Physics, \textbf{2021} (2021), no.~3, 178,
\href{https://doi.org/10.1007/JHEP03(2021)178}{DOI: 10.1007/JHEP03(2021)178}.

\bibitem{EkelandHofer1989}
I.~Ekeland and H.~Hofer,
\emph{Symplectic topology and Hamiltonian dynamics},
Mathematische Zeitschrift, \textbf{200} (1989), no.~3, 355--378,
\href{https://doi.org/10.1007/BF01215653}{DOI: 10.1007/BF01215653}.

\bibitem{Federer1959}
H.~Federer,
\emph{Curvature measures},
Transactions of the American Mathematical Society, \textbf{93} (1959), no.~3, 418--491,
\href{https://doi.org/10.1090/S0002-9947-1959-0110078-1}{DOI: 10.1090/S0002-9947-1959-0110078-1}.

\bibitem{GibsonEtAl1976}
C.~G. Gibson, K.~Wirthm\"uller, A.~A. du~Plessis, and E.~J.~N. Looijenga,
\emph{Topological Stability of Smooth Mappings},
Lecture Notes in Mathematics, vol.~552, Springer, Berlin, Heidelberg, 1976,
\href{https://doi.org/10.1007/BFb0095244}{DOI: 10.1007/BFb0095244}.

\bibitem{Gromov1987}
M.~Gromov,
\emph{Hyperbolic groups},
in \emph{Essays in Group Theory}, ed. S.~M. Gersten, Mathematical Sciences Research Institute Publications, vol.~8, Springer, New York, 1987, pp.~75--263,
ISBN: 978-0-387-96618-2,
\href{https://doi.org/10.1007/978-1-4613-9586-7_3}{DOI: 10.1007/978-1-4613-9586-7\_3}.

\bibitem{HoferZehnder1994}
H.~Hofer and E.~Zehnder,
\emph{Symplectic Invariants and Hamiltonian Dynamics},
Birkh\"auser Advanced Texts: Basler Lehrb\"ucher, Birkh\"auser Basel, 1994,
ISBN: 978-3-0348-8540-9,
\href{https://doi.org/10.1007/978-3-0348-8540-9}{DOI: 10.1007/978-3-0348-8540-9}.

\bibitem{Jarzynski1997}
C.~Jarzynski,
\emph{Nonequilibrium equality for free energy differences},
Physical Review Letters, \textbf{78} (1997), no.~14, 2690--2693,
\href{https://doi.org/10.1103/PhysRevLett.78.2690}{DOI: 10.1103/PhysRevLett.78.2690}.

\bibitem{KashiwaraSchapira1990}
M.~Kashiwara and P.~Schapira,
\emph{Sheaves on Manifolds},
Grundlehren der mathematischen Wissenschaften, vol.~292, Springer-Verlag, Berlin, Heidelberg, 1990,
ISBN: 978-3-540-51861-7,
\href{https://doi.org/10.1007/978-3-662-02661-8}{DOI: 10.1007/978-3-662-02661-8}.

\bibitem{BridsonHaefliger1999}
M.~R. Bridson and A.~Haefliger,
\emph{Metric Spaces of Non-Positive Curvature},
Grundlehren der mathematischen Wissenschaften, vol.~319, Springer-Verlag, Berlin, Heidelberg, 1999,
ISBN: 978-3-540-64324-1,
\href{https://doi.org/10.1007/978-3-662-12494-9}{DOI: 10.1007/978-3-662-12494-9}.

\bibitem{JuutinenLindqvistManfredi1999}
P.~Juutinen, P.~Lindqvist, and J.~J. Manfredi,
\emph{The $\infty$-eigenvalue problem},
Archive for Rational Mechanics and Analysis, \textbf{148} (1999), no.~2, 89--105,
\href{https://doi.org/10.1007/s002050050157}{DOI: 10.1007/s002050050157}.

\bibitem{OttoVillani2000}
F.~Otto and C.~Villani,
\emph{Generalization of an inequality by Talagrand and links with the logarithmic Sobolev inequality},
Journal of Functional Analysis, \textbf{173} (2000), no.~2, 361--400,
\href{https://doi.org/10.1006/jfan.1999.3557}{DOI: 10.1006/jfan.1999.3557}.

\bibitem{KawohlFridman2003}
B.~Kawohl and V.~Fridman,
\emph{Isoperimetric estimates for the first eigenvalue of the $p$-Laplace operator and the Cheeger constant},
Commentationes Mathematicae Universitatis Carolinae, \textbf{44} (2003), no.~4, 659--667,
\href{https://eudml.org/doc/249339}{EuDML: 249339}, \href{https://dml.cz/handle/10338.dmlcz/119420}{DML-CZ: 119420}.

\bibitem{Lindqvist1990}
P.~Lindqvist,
\emph{On the equation $\operatorname{div}(|\nabla u|^{p-2}\nabla u) + \lambda|u|^{p-2}u = 0$},
Proceedings of the American Mathematical Society, \textbf{109} (1990), no.~1, 157--164,
\href{https://doi.org/10.1090/S0002-9939-1990-1007505-7}{DOI: 10.1090/S0002-9939-1990-1007505-7}.

\bibitem{SilvaFilhoBTS}
R.~M. Silva-Filho,
\emph{Beyond the Spectrum: The Complete Three-Volume Monograph on Functional Tensor Realizations, Metric Measure Geometry, and Higher Topological Invariants},
Zenodo (2026),
\href{https://doi.org/10.5281/zenodo.22644743}{DOI: 10.5281/zenodo.22644743}.

\bibitem{SivakCrooks2012}
D.~A. Sivak and G.~E. Crooks,
\emph{Thermodynamic metrics and optimal paths},
Physical Review Letters, \textbf{108} (2012), no.~19, 190602,
\href{https://doi.org/10.1103/PhysRevLett.108.190602}{DOI: 10.1103/PhysRevLett.108.190602}.

\bibitem{ArtsteinAvidanKarasevOstrover2014}
S.~Artstein-Avidan, R.~Karasev, and Y.~Ostrover,
\emph{From symplectic measurements to the Mahler conjecture},
Duke Mathematical Journal, \textbf{163} (2014), no.~11,
\href{https://doi.org/10.1215/00127094-2794999}{DOI: 10.1215/00127094-2794999}.

\bibitem{ArtsteinAvidanOstrover2008}
S.~Artstein-Avidan and Y.~Ostrover,
\emph{A Brunn--Minkowski inequality for symplectic capacities of convex domains},
International Mathematics Research Notices IMRN (2008), rnn044,
\href{https://doi.org/10.1093/imrn/rnn044}{DOI: 10.1093/imrn/rnn044}.

\bibitem{AkersRath2020}
C.~Akers and P.~Rath,
\emph{Entanglement wedge cross sections require tripartite entanglement},
Journal of High Energy Physics, \textbf{2020} (2020), no.~4, 208,
\href{https://doi.org/10.1007/JHEP04(2020)208}{DOI: 10.1007/JHEP04(2020)208}.

\bibitem{BelloniKawohlJuutinen2006}
M.~Belloni, B.~Kawohl, and P.~Juutinen,
\emph{The $p$-Laplace eigenvalue problem as $p\to\infty$ in a Finsler metric},
Journal of the European Mathematical Society, \textbf{8} (2006), no.~1, 123--138,
\href{https://doi.org/10.4171/JEMS/40}{DOI: 10.4171/JEMS/40}.

\bibitem{BlaberSivak2023}
S.~Blaber and D.~A. Sivak,
\emph{Optimal control in stochastic thermodynamics},
Journal of Physics Communications, \textbf{7} (2023), no.~3, 033001,
\href{https://doi.org/10.1088/2399-6528/acbf04}{DOI: 10.1088/2399-6528/acbf04}.

\bibitem{BoissonnatLieutierWintraecken2019}
J.-D. Boissonnat, A.~Lieutier, and M.~Wintraecken,
\emph{The reach, metric distortion, geodesic convexity and the variation of tangent spaces},
Journal of Applied and Computational Topology, \textbf{3} (2019), 29--58,
\href{https://doi.org/10.1007/s41468-019-00029-8}{DOI: 10.1007/s41468-019-00029-8}.

\bibitem{Carlson1963}
B.~C. Carlson,
\emph{Lauricella's hypergeometric function $F_D$},
Journal of Mathematical Analysis and Applications, \textbf{7} (1963), no.~3, 452--470,
\href{https://doi.org/10.1016/0022-247X(63)90067-2}{DOI: 10.1016/0022-247X(63)90067-2}.

\bibitem{ChampionDePascaleJimenez2008}
T.~Champion, L.~De~Pascale, and C.~Jimenez,
\emph{The $\infty$-eigenvalue problem and a problem of optimal transportation},
preprint (2008), arXiv:0811.1934,
\href{https://doi.org/10.48550/arXiv.0811.1934}{DOI: 10.48550/arXiv.0811.1934}.

\bibitem{Cheeger1970}
J.~Cheeger,
\emph{A lower bound for the smallest eigenvalue of the Laplacian},
in \emph{Problems in Analysis: A Symposium in Honor of Salomon Bochner}, Princeton University Press, Princeton, 1970, pp.~195--200,
\href{https://doi.org/10.1515/9781400869312-013}{DOI: 10.1515/9781400869312-013}.

\bibitem{Chow2022}
D.~D.~K. Chow,
\emph{Schl\"omilch integrals and probability distributions on the simplex},
preprint (2022), arXiv:2201.11013,
\href{https://doi.org/10.48550/arXiv.2201.11013}{DOI: 10.48550/arXiv.2201.11013}.

\bibitem{CieliebakHoferLatschevSchlenk2005}
K.~Cieliebak, H.~Hofer, J.~Latschev, and F.~Schlenk,
\emph{Quantitative symplectic geometry},
preprint (2005), arXiv:math/0506191,
\href{https://doi.org/10.48550/arXiv.math/0506191}{DOI: 10.48550/arXiv.math/0506191}.

\bibitem{Crooks2007}
G.~E. Crooks,
\emph{Measuring thermodynamic length},
Physical Review Letters, \textbf{99} (2007), no.~10, 100602,
\href{https://doi.org/10.1103/PhysRevLett.99.100602}{DOI: 10.1103/PhysRevLett.99.100602}.

\bibitem{DLMF}
F.~W.~J. Olver, D.~W. Lozier, R.~F. Boisvert, and C.~W. Clark (eds.),
\emph{NIST Handbook of Mathematical Functions},
Cambridge University Press, Cambridge, 2010, ISBN: 978-0-521-19225-5;
online as the \emph{NIST Digital Library of Mathematical Functions}, \url{https://dlmf.nist.gov}.

\bibitem{EkelandHofer1990}
I.~Ekeland and H.~Hofer,
\emph{Symplectic topology and Hamiltonian dynamics~II},
Mathematische Zeitschrift, \textbf{203} (1990), no.~1, 553--567,
\href{https://doi.org/10.1007/BF02570756}{DOI: 10.1007/BF02570756}.

\bibitem{FloerHofer1994}
A.~Floer and H.~Hofer,
\emph{Symplectic homology I: Open sets in $\mathbb{C}^n$},
Mathematische Zeitschrift, \textbf{215} (1994), no.~1, 37--88,
\href{https://doi.org/10.1007/BF02571699}{DOI: 10.1007/BF02571699}.

\bibitem{FukagaiItoNarukawa1999}
N.~Fukagai, M.~Ito, and K.~Narukawa,
\emph{Limit as $p\to\infty$ of $p$-Laplace eigenvalue problems and $L^\infty$-inequality of the Poincar\'e type},
Differential and Integral Equations, \textbf{12} (1999), no.~2, 183--206,
\href{https://doi.org/10.57262/die/1367265629}{DOI: 10.57262/die/1367265629}.

\bibitem{Gromov1985}
M.~Gromov,
\emph{Pseudo holomorphic curves in symplectic manifolds},
Inventiones Mathematicae, \textbf{82} (1985), no.~2, 307--347,
\href{https://doi.org/10.1007/BF01388806}{DOI: 10.1007/BF01388806}.

\bibitem{HaimKislevOstrover2026}
P.~Haim-Kislev and Y.~Ostrover,
\emph{A counterexample to Viterbo's conjecture},
Annals of Mathematics, \textbf{203} (2026), no.~2, 603--622,
\href{https://doi.org/10.4007/annals.2026.203.2.5}{DOI: 10.4007/annals.2026.203.2.5}.

\bibitem{HaydenLemmSorce2023}
P.~Hayden, M.~Lemm, and J.~Sorce,
\emph{Reflected entropy is not a correlation measure},
Physical Review A, \textbf{107} (2023), no.~5, L050401,
\href{https://doi.org/10.1103/PhysRevA.107.L050401}{DOI: 10.1103/PhysRevA.107.L050401}.

\bibitem{HaydenParrikarSorce2021}
P.~Hayden, O.~Parrikar, and J.~Sorce,
\emph{The Markov gap for geometric reflected entropy},
Journal of High Energy Physics, \textbf{2021} (2021), no.~10, 47,
\href{https://doi.org/10.1007/JHEP10(2021)047}{DOI: 10.1007/JHEP10(2021)047}.

\bibitem{JordanKinderlehrerOtto1998}
R.~Jordan, D.~Kinderlehrer, and F.~Otto,
\emph{The variational formulation of the Fokker--Planck equation},
SIAM Journal on Mathematical Analysis, \textbf{29} (1998), no.~1, 1--17,
\href{https://doi.org/10.1137/S0036141096303359}{DOI: 10.1137/S0036141096303359}.

\bibitem{Kashiwara1985}
M.~Kashiwara,
\emph{Index theorem for constructible sheaves},
Ast\'erisque, \textbf{130} (1985), 193--209,
\href{http://www.numdam.org/item?id=AST_1985__130__193_0}{Numdam: AST\_1985\_\_130\_\_193\_0}.

\bibitem{LiebRuskai1973}
E.~H. Lieb and M.~B. Ruskai,
\emph{Proof of the strong subadditivity of quantum-mechanical entropy},
Journal of Mathematical Physics, \textbf{14} (1973), no.~12, 1938--1941,
\href{https://doi.org/10.1063/1.1666274}{DOI: 10.1063/1.1666274}.

\bibitem{Lindqvist2008}
P.~Lindqvist,
\emph{A nonlinear eigenvalue problem},
in \emph{Topics in Mathematical Analysis}, Series on Analysis, Applications and Computation, World Scientific, Singapore, 2008, pp.~175--203,
ISBN: 978-981-281-105-9,
\href{https://doi.org/10.1142/9789812811066_0005}{DOI: 10.1142/9789812811066\_0005}.

\bibitem{MacPherson1974}
R.~D. MacPherson,
\emph{Chern classes for singular algebraic varieties},
Annals of Mathematics, \textbf{100} (1974), no.~2, 423--432,
\href{https://doi.org/10.2307/1971080}{DOI: 10.2307/1971080}.

\bibitem{Mather2012}
J.~Mather,
\emph{Notes on topological stability},
Bulletin of the American Mathematical Society, \textbf{49} (2012), no.~4, 475--506,
\href{https://doi.org/10.1090/S0273-0979-2012-01383-6}{DOI: 10.1090/S0273-0979-2012-01383-6}.

\bibitem{Mercer1909}
J.~Mercer,
\emph{Functions of positive and negative type, and their connection with the theory of integral equations},
Philosophical Transactions of the Royal Society of London, Series~A, \textbf{209} (1909), 415--446,
\href{https://doi.org/10.1098/rsta.1909.0016}{DOI: 10.1098/rsta.1909.0016}.

\bibitem{NakazatoIto2021}
M.~Nakazato and S.~Ito,
\emph{Geometrical aspects of entropy production in stochastic thermodynamics based on Wasserstein distance},
Physical Review Research, \textbf{3} (2021), no.~4, 043093,
\href{https://doi.org/10.1103/PhysRevResearch.3.043093}{DOI: 10.1103/PhysRevResearch.3.043093}.

\bibitem{NguyenEtAl2018}
P.~Nguyen, T.~Devakul, M.~G. Halbasch, M.~P. Zaletel, and B.~Swingle,
\emph{Entanglement of purification: from spin chains to holography},
Journal of High Energy Physics, \textbf{2018} (2018), no.~1, 98,
\href{https://doi.org/10.1007/JHEP01(2018)098}{DOI: 10.1007/JHEP01(2018)098}.

\bibitem{NiyogiSmaleWeinberger2008}
P.~Niyogi, S.~Smale, and S.~Weinberger,
\emph{Finding the homology of submanifolds with high confidence from random samples},
Discrete \& Computational Geometry, \textbf{39} (2008), 419--441,
\href{https://doi.org/10.1007/s00454-008-9053-2}{DOI: 10.1007/s00454-008-9053-2}.

\bibitem{OrroTrotman2010}
P.~Orro and D.~Trotman,
\emph{Regularity of the transverse intersection of two regular stratifications},
in \emph{Real and Complex Singularities}, London Mathematical Society Lecture Note Series, vol.~380, Cambridge University Press, Cambridge, 2010, pp.~298--304,
ISBN: 978-0-521-16969-1,
\href{https://doi.org/10.1017/CBO9780511731983.022}{DOI: 10.1017/CBO9780511731983.022}.

\bibitem{RatajZahle2019}
J.~Rataj and M.~Z\"ahle,
\emph{Curvature Measures of Singular Sets},
Springer Monographs in Mathematics, Springer, Cham, 2019,
ISBN: 978-3-030-18182-6,
\href{https://doi.org/10.1007/978-3-030-18183-3}{DOI: 10.1007/978-3-030-18183-3}.

\bibitem{SalamonBerry1983}
P.~Salamon and R.~S. Berry,
\emph{Thermodynamic length and dissipated availability},
Physical Review Letters, \textbf{51} (1983), no.~13, 1127--1130,
\href{https://doi.org/10.1103/PhysRevLett.51.1127}{DOI: 10.1103/PhysRevLett.51.1127}.

\bibitem{Schapira1991}
P.~Schapira,
\emph{Operations on constructible functions},
Journal of Pure and Applied Algebra, \textbf{72} (1991), no.~1, 83--93,
\href{https://doi.org/10.1016/0022-4049(91)90131-K}{DOI: 10.1016/0022-4049(91)90131-K}.

\bibitem{Schapira2017}
P.~Schapira,
\emph{Microlocal analysis and beyond},
preprint, 2017, DOI:
\href{https://doi.org/10.48550/arXiv.1701.08955}{10.48550/arXiv.1701.08955}.

\bibitem{SchmiedlSeifert2007}
T.~Schmiedl and U.~Seifert,
\emph{Optimal finite-time processes in stochastic thermodynamics},
Physical Review Letters, \textbf{98} (2007), no.~10, 108301,
\href{https://doi.org/10.1103/PhysRevLett.98.108301}{DOI: 10.1103/PhysRevLett.98.108301}.

\bibitem{Schmidt1907}
E.~Schmidt,
\emph{Zur Theorie der linearen und nichtlinearen Integralgleichungen. I.~Teil},
Mathematische Annalen, \textbf{63} (1907), no.~4, 433--476,
\href{https://doi.org/10.1007/BF01449770}{DOI: 10.1007/BF01449770}.

\bibitem{Schwarz2000}
M.~Schwarz,
\emph{On the action spectrum for closed symplectically aspherical manifolds},
Pacific Journal of Mathematics, \textbf{193} (2000), no.~2, 419--461,
\href{https://doi.org/10.2140/pjm.2000.193.419}{DOI: 10.2140/pjm.2000.193.419}.

\bibitem{Seifert2012}
U.~Seifert,
\emph{Stochastic thermodynamics, fluctuation theorems and molecular machines},
Reports on Progress in Physics, \textbf{75} (2012), no.~12, 126001,
\href{https://doi.org/10.1088/0034-4885/75/12/126001}{DOI: 10.1088/0034-4885/75/12/126001}.

\bibitem{silvafilho2026book}
R.~M. Silva-Filho,
\emph{Geometry, Tensors, and Quantum Gravity},
monograph, Zenodo (2026),
\href{https://doi.org/10.5281/zenodo.22290043}{DOI: 10.5281/zenodo.22290043}.

\bibitem{Srednicki1993}
M.~Srednicki,
\emph{Entropy and area},
Physical Review Letters, \textbf{71} (1993), no.~5, 666--669,
\href{https://doi.org/10.1103/PhysRevLett.71.666}{DOI: 10.1103/PhysRevLett.71.666}.

\bibitem{Talagrand1996}
M.~Talagrand,
\emph{Transportation cost for Gaussian and other product measures},
Geometric and Functional Analysis, \textbf{6} (1996), no.~3, 587--600,
\href{https://doi.org/10.1007/BF02249265}{DOI: 10.1007/BF02249265}.

\bibitem{Trotman2020}
D.~Trotman,
\emph{Stratification theory},
in \emph{Handbook of Geometry and Topology of Singularities~I}, Springer, Cham, 2020, pp.~243--273,
ISBN: 978-3-030-53060-0,
\href{https://doi.org/10.1007/978-3-030-53061-7_4}{DOI: 10.1007/978-3-030-53061-7\_4}.

\bibitem{UmemotoTakayanagi2018}
K.~Umemoto and T.~Takayanagi,
\emph{Entanglement of purification through holographic duality},
Nature Physics, \textbf{14} (2018), no.~6, 573--577,
\href{https://doi.org/10.1038/s41567-018-0075-2}{DOI: 10.1038/s41567-018-0075-2}.

\bibitem{VanVuSaito2023}
T.~Van~Vu and K.~Saito,
\emph{Thermodynamic unification of optimal transport: thermodynamic uncertainty relation, minimum dissipation, and thermodynamic speed limits},
Physical Review X, \textbf{13} (2023), no.~1, 011013,
\href{https://doi.org/10.1103/PhysRevX.13.011013}{DOI: 10.1103/PhysRevX.13.011013}.

\bibitem{Villani2009}
C.~Villani,
\emph{Optimal Transport: Old and New},
Grundlehren der mathematischen Wissenschaften, vol.~338, Springer, Berlin, Heidelberg, 2009,
ISBN: 978-3-540-71049-3,
\href{https://doi.org/10.1007/978-3-540-71050-9}{DOI: 10.1007/978-3-540-71050-9}.

\bibitem{Viro1988}
O.~Y. Viro,
\emph{Some integral calculus based on Euler characteristic},
in \emph{Topology and Geometry --- Rohlin Seminar}, Lecture Notes in Mathematics, vol.~1346, Springer, Berlin, Heidelberg, 1988, pp.~127--138,
ISBN: 978-3-540-50237-1,
\href{https://doi.org/10.1007/BFb0082775}{DOI: 10.1007/BFb0082775}.

\bibitem{Viterbo1992}
C.~Viterbo,
\emph{Symplectic topology as the geometry of generating functions},
Mathematische Annalen, \textbf{292} (1992), no.~1, 685--710,
\href{https://doi.org/10.1007/BF01444643}{DOI: 10.1007/BF01444643}.

\bibitem{Viterbo2000}
C.~Viterbo,
\emph{Metric and isoperimetric problems in symplectic geometry},
Journal of the American Mathematical Society, \textbf{13} (2000), no.~2, 411--431,
\href{https://doi.org/10.1090/S0894-0347-00-00328-3}{DOI: 10.1090/S0894-0347-00-00328-3}.

\bibitem{vonRenesseSturm2005}
M.-K. von~Renesse and K.-T. Sturm,
\emph{Transport inequalities, gradient estimates, entropy and Ricci curvature},
Communications on Pure and Applied Mathematics, \textbf{58} (2005), no.~7, 923--940,
\href{https://doi.org/10.1002/cpa.20060}{DOI: 10.1002/cpa.20060}.

\bibitem{Weyl1939}
H.~Weyl,
\emph{On the volume of tubes},
American Journal of Mathematics, \textbf{61} (1939), no.~2, 461--472,
\href{https://doi.org/10.2307/2371513}{DOI: 10.2307/2371513}.

\bibitem{Whitney1965}
H.~Whitney,
\emph{Tangents to an analytic variety},
Annals of Mathematics, \textbf{81} (1965), no.~3, 496--549,
\href{https://doi.org/10.2307/1970400}{DOI: 10.2307/1970400}.

\bibitem{ZhongDeWeese2024}
A.~Zhong and M.~R. DeWeese,
\emph{Beyond linear response: equivalence between thermodynamic geometry and optimal transport},
Physical Review Letters, \textbf{133} (2024), no.~5, 057102,
\href{https://doi.org/10.1103/PhysRevLett.133.057102}{DOI: 10.1103/PhysRevLett.133.057102}.

\end{thebibliography}

\end{document}
